\documentclass[11pt,letterpaper,reqno]{amsart}
% BEGIN INLINE manuscript-style
\makeatletter
\NeedsTeXFormat{LaTeX2e}[2020/10/01]
% Use with \documentclass[11pt,letterpaper,reqno]{amsart}.
% This file controls presentation only. It supplies no mathematical content.

\RequirePackage[T1]{fontenc}
\RequirePackage{lmodern}
\RequirePackage{microtype}
\RequirePackage[letterpaper,margin=1in,headheight=15pt]{geometry}
\RequirePackage{amsmath,amsthm,amssymb,mathtools}
\RequirePackage{mathrsfs}
\RequirePackage{enumitem}
\RequirePackage{perpage}
\RequirePackage{longtable}
\RequirePackage{xcolor}
\definecolor{manuscriptred}{rgb}{0.6,0.1,0.1}
\definecolor{manuscriptblue}{rgb}{0.1,0.3,0.5}
\definecolor{manuscriptmagenta}{HTML}{60445C}
\RequirePackage[breaklinks=true,colorlinks=true,
  linkcolor=manuscriptred,citecolor=manuscriptblue,urlcolor=manuscriptmagenta]{hyperref}
\RequirePackage{cleveref}
\hypersetup{pdfauthor={}}

% Retain the AMS paragraph, title, heading and running-head conventions.
% Equation counters restart each section: 2.i, 2.ii, ... .
% Appendix sections keep their letters: A.i, A.ii, ... .
\numberwithin{equation}{section}
\renewcommand{\theequation}{\thesection.\roman{equation}}
% Prose references include "eq." and omit parentheses around the number.
\renewcommand{\eqref}[1]{\textup{eq.~\ref{#1}}}
\crefname{equation}{eq.}{eqs.}
\Crefname{equation}{eq.}{eqs.}
\creflabelformat{equation}{#2#1#3}
\setcounter{tocdepth}{1}
\renewcommand{\thefigure}{\Roman{figure}}
\renewcommand{\thetable}{\Roman{table}}
\newcommand{\fpneedspace}[1]{\par\begingroup
  \dimen0=#1\relax
  \vskip0pt plus\dimen0\penalty-100
  \vskip0pt plus-\dimen0\vskip\dimen0
  \penalty9999\vskip-\dimen0\vskip0pt\endgroup}
\MakePerPage{footnote}
\renewcommand{\thefootnote}{%
  \ifnum\value{footnote}>9\arabic{footnote}\else\fnsymbol{footnote}\fi}

% Upright statements. Main results are A, B, C; supporting statements
% share a section-local sequence. These end marks indicate typography,
% never proof completeness, validation, or a deferred-proof status.
\theoremstyle{definition}
\newtheorem{ms@theorem}{Theorem}
\renewcommand{\thems@theorem}{\Alph{ms@theorem}}
\newtheorem*{ms@theoremstar}{Theorem}
\newtheorem*{ms@theoremzero}{Theorem 0}
\newtheorem{ms@proposition}{Proposition}[section]
\newtheorem{ms@lemma}[ms@proposition]{Lemma}
\newtheorem{ms@definition}[ms@proposition]{Definition}
\newtheorem{ms@corollary}[ms@proposition]{Corollary}
\newtheorem*{ms@corollarystar}{Corollary}

\newcommand{\ms@majorstart}{%
  \pushQED{\qed}\renewcommand{\qedsymbol}{$\blacksquare\square$}}
\newcommand{\ms@minorstart}{%
  \pushQED{\qed}\renewcommand{\qedsymbol}{$\blacksquare$}}
\newenvironment{theorem}[1][]{\ms@majorstart\ms@theorem[#1]}
  {\popQED\endms@theorem}
\newenvironment{theorem*}[1][]{\ms@majorstart\ms@theoremstar[#1]}
  {\popQED\endms@theoremstar}
\newenvironment{theoremzero}[1][]{\ms@majorstart\ms@theoremzero[#1]}
  {\popQED\endms@theoremzero}
\newenvironment{proposition}[1][]{\ms@minorstart\ms@proposition[#1]}
  {\popQED\endms@proposition}
\newenvironment{lemma}[1][]{\ms@minorstart\ms@lemma[#1]}
  {\popQED\endms@lemma}
\newenvironment{definition}[1][]{\ms@minorstart\ms@definition[#1]}
  {\popQED\endms@definition}
\newenvironment{corollary}[1][]{\ms@majorstart\ms@corollary[#1]}
  {\popQED\endms@corollary}
\newenvironment{corollary*}[1][]{\ms@majorstart\ms@corollarystar[#1]}
  {\popQED\endms@corollarystar}
% The source's lemma/proposition variants are deliberately not carried over:
% choose one marker for each statement type rather than encode proof status.

\theoremstyle{remark}
\newtheorem{remark}[ms@proposition]{Remark}
\newtheorem*{remark*}{Remark}
% The ordinary amsthm proof environment retains its open square.

\crefname{ms@theorem}{theorem}{theorems}
\Crefname{ms@theorem}{Theorem}{Theorems}
\crefname{ms@proposition}{proposition}{propositions}
\Crefname{ms@proposition}{Proposition}{Propositions}
\crefname{ms@lemma}{lemma}{lemmas}
\Crefname{ms@lemma}{Lemma}{Lemmas}
\crefname{ms@definition}{definition}{definitions}
\Crefname{ms@definition}{Definition}{Definitions}
\crefname{ms@corollary}{corollary}{corollaries}
\Crefname{ms@corollary}{Corollary}{Corollaries}
\crefname{remark}{remark}{remarks}
\Crefname{remark}{Remark}{Remarks}

% Bibliography configuration for this manuscript.
\makeatother
% END INLINE manuscript-style

\title[Physical short-distance asymptotics]{Physical short-distance asymptotics of the lattice Yang--Mills action density}
\date{Working draft, 9 November 2026}
\begin{document}
\raggedbottom
\begin{abstract}
We derive a two-point action-density comparison at a prescribed
short-distance scale from explicit local comparison and running
coefficient estimates. The source factor is the derivative of the
actual running coupling shift with the real reference data fixed.
The normalized free curvature covariance gives the coefficient
\(9/\pi^4\) for \(SU(2)\). Exact source differentiation, local
response estimates and a conditional reference-law comparison
control the physical errors relative to the leading fourth power
of the running coupling. A common holomorphic cell normalizer
transfers the comparison to every admitted periodic Gibbs state
at the same bare datum. The result gives a logarithmically
decreasing ultraviolet upper bound and a positive lower bound
for separated nonnegative tests. The actual local conditional
law, history estimates and running coefficient bounds remain
explicit hypotheses; continuum existence is not asserted.
\end{abstract}
\maketitle
\tableofcontents

\section{Introduction}

The short-distance covariance of the Wilson action density has
two distinct sources of normalization. The quadratic curvature
law fixes its free coefficient. The response of the extracted
couplings to a scalar source fixes the factor multiplying each
physical insertion. Both must be identified before a comparison
with a Gaussian measure has a definite meaning. A bound that is
merely small in absolute terms is insufficient when the leading
covariance itself tends to zero with the running coupling.

We first fix the physical observable and its free normalization.
The running coefficient estimates then select the observation
depth and determine the constant-source response with all real
reference data held fixed. Differentiation of the actual retained
source weight separates the covariance of its first logarithmic
jets from its mixed logarithmic jet. The local comparison hypotheses control
these terms, including the contributions of retained histories
and the conditional exterior field.

The passage from a reference torus to arbitrary periodic Gibbs
states is a separate step. A common cell normalizer permits
source derivatives to be compared without estimating their
possibly divergent individual means. Signed copies of a test
pair improve the resulting state difference by the number of
disjoint copies. The error is then negligible relative to the
positive free signal on the full prescribed shape class.

The extraction of a scalar curl coefficient from an effective
action follows the polarization analysis of
Ba\l aban \cite[\S5, pp.~292--297]{Bal87}. That work treats a
small-field approximation. It does not provide the physical
conditional-law and history estimates assumed here. We use the
free curvature calculation of \cite{GI}, the admitted trajectory
and normalization conventions of \cite{RT}, and the fixed-source
and state-transfer results of \cite{SI}, with reflection and
chessboard estimates from \cite{LP}. The hypotheses below state
the additional local comparison and coefficient information
needed for the physical law. Their establishment and continuum
reconstruction are separate problems.

\section{Physical observables and fixed source data}
\label{fp:sec-setting}

Fix an odd integer block factor \(L\ge3\). At cutoff depth
\(K\in\mathbb N\), the mesh is \(a=L^{-K}\). A compatible
torus is \(\Lambda_{a,\ell}=(a\mathbb Z/\ell\mathbb Z)^4\),
with \(\ell/a\) integral and with the block divisibilities required
at the depths in use. One independent bond variable lies in
\(SU(2)\) for each positively oriented nearest-neighbour bond;
reversing a bond replaces its variable by its inverse. Haar measure
has mass one. For a positively oriented plaquette \(p\), let
\(U_p\) be its ordered bond product and \(x_p\) its barycentre.
Put
\begin{equation}\label{fp:eq-physical-observable}
 A(U)=\sum_p a_p(U),\qquad
 a_p(U)=1-\tfrac12\operatorname{tr}U_p,\qquad
 \mathcal O_a(f)=\sum_p f(x_p)a_p(U).
\end{equation}
Here \(f\) is a compactly supported real Lipschitz function, or
its periodic version on the torus. The continuum test class will
be the \(C_c^1\) subclass. The six positive plaquette
orientations are all included. The trace is real on \(SU(2)\),
and \(0\le a_p\le2\). In particular, every fixed-cutoff compact
insertion is bounded.

The torus Wilson measure is the normalized measure
\(e^{-\beta A(U)}dU\), where \(\beta\ge0\). On the infinite
lattice, an admitted state is a periodic Wilson Gibbs state at
this same pair \((a,\beta)\). Periodicity means invariance under
a finite-index subgroup of bare lattice translations. The
reflection theorem of \cite[Theorem B]{LP} gives invariance under
all bare lattice translations for these states. It does not
identify every such state with a particular sequence of torus
measures. For any of these probability measures \(\rho\), write
\begin{equation}\label{fp:eq-physical-covariance}
 S^T_{2;a,\rho}(f,h)
   =\rho\bigl(\mathcal O_a(f)\mathcal O_a(h)\bigr)
      -\rho\bigl(\mathcal O_a(f)\bigr)
       \rho\bigl(\mathcal O_a(h)\bigr).
\end{equation}
This is a bilinear covariance, without complex conjugation.
Complex sources are used to differentiate it; the comparison
and positivity conclusions concern real nonnegative tests.

There is no additional \(a^4\) in
\eqref{fp:eq-physical-observable}. To check the convention, use
traceless Hermitian colour coordinates with
\(\langle X,Y\rangle=\tfrac12\operatorname{tr}(XY)\).
For a smooth field,
\(U_p=\exp(i a^2F_{\mu\nu}(x_p)+O(a^3))\) implies
\(a_p=\tfrac12a^4|F_{\mu\nu}(x_p)|^2+O(a^5)\).
The plaquette sum therefore approximates
\(\tfrac12\int f\sum_{\mu<\nu}|F_{\mu\nu}|^2\).
This expansion fixes the normalization of the free quadratic
density. It is not a smoothness assertion about the random field.

For finitely many tests, let
\(W_z=\sum_i z_i f_i\). On a finite torus set
\[
 Z_{a,\ell}(W_z)=\int
       \exp\{-\beta A(U)+\mathcal O_a(W_z)\}\,dU.
\]
For an infinite-volume state use
\(\rho(e^{\mathcal O_a(W_z)})\) instead of the partition ratio.
Both functions are entire in the finite source vector at fixed
cutoff: the exponent is a finite linear combination of bounded
observables, so its power series is uniformly dominated on every
compact source set. Both have value one after normalization at
zero, and hence have a local holomorphic logarithm there.

All real reference couplings, test cutoffs, averaging kernels,
coordinate choices and history labels are fixed before the
source is varied. In an exact blocked representation, its
positive reference measure is likewise the unsourced measure.
The source dependence is carried by the actual retained weight. The
transport identity of \cite[Proposition 3.1]{SI} supplies this
representation when its fixed-kernel hypotheses hold. A new
reference construction at each source value would require the
derivatives of that construction as well and is not used here.

The running inverse coupling at a witness depth will be denoted
by \(x=x_{K}(d)\), with \(g_K(d)=x^{-1/2}\). Its constant-source
response is determined from the actual coefficient maps below.
The fixed small terminal coupling that selects the bare sequence
is a different parameter. Constants may depend on that terminal
choice, \(L\), the fixed construction parameters and the fixed
shape threshold. They are independent of cutoff depth, admitted
state, test position and test amplitudes unless a statement says
otherwise. Uniform error functions are always stated on the
admitted parameter range. All logarithms are natural.


\section{The separated free comparison}
\label{fp:sec-free}

\subsection{Test geometry}

Write \(T_\ell=\mathbb R^4/\ell\mathbb Z^4\), with its flat distance and
unnormalized volume measure. Let \(Y=\mathbb R^4\) or \(T_\ell\).
\begin{definition}[Separated test pairs]\label{fp:def-shape}
For \(d>0\) and \(0<\vartheta\le1\), let
\(\mathfrak T_Y(d;\vartheta)\) consist of pairs \((f,h)\) of nonzero,
nonnegative Lipschitz functions. Each support lies in a ball of
radius \(d\). Write \(M_v=\|v\|_\infty\) and \(m_v=\int_Y v\)
for \(v=f,h\).
The requirements are
\begin{equation}\label{fp:eq-shape}
 \begin{gathered}
 d\le\operatorname{dist}(\operatorname{supp}f,\operatorname{supp}h)
       \le6d,\\
 m_v\ge\vartheta M_vd^4,\qquad
 \operatorname{Lip}(v)\le M_v/(\vartheta d).
 \end{gathered}
\end{equation}
On \(T_\ell\) take \(\ell\ge40d\).
The \(C_c^1\) class used in the continuum application is contained in
this class with the same parameter \(\vartheta_*=1/12000\).
\end{definition}

\begin{lemma}[Size and scaling of the comparison integral]
\label{fp:lem-shape-kernel}
For \((f,h)\in\mathfrak T_Y(d;\vartheta)\), the two supports lie in a
common ball of radius \(5d\). On their product the distance lies in
\([d,10d]\). With \(v_4=\pi^2/2\), the integral
\[
 \mathcal K(f,h)=\iint f(x)h(y)|x-y|^{-8}\,dx\,dy
\]
satisfies
\begin{equation}\label{fp:eq-shape-kernel-bounds}
 10^{-8}\vartheta^2 M_fM_h
 \le (10d)^{-8}m_fm_h
 \le\mathcal K(f,h)
 \le d^{-8}m_fm_h
 \le v_4^2 M_fM_h .
\end{equation}
On a torus the integral is evaluated using a common lift of the supports
to \(\mathbb R^4\). It is independent of the lift. For \(t>0\) let
\(v_t(x)=v(x/t)\) on \(tY\). Then
\((f_t,h_t)\in\mathfrak T_{tY}(td;\vartheta)\) and
\(\mathcal K(f_t,h_t)=\mathcal K(f,h)\).
\end{lemma}

\begin{proof}
Choose points attaining the minimum cross-support distance. The
distance of the two supporting ball centres is at most \(d+6d+d=8d\).
Choose lifts of those centres realizing that distance when \(Y\)
is a torus. Their midpoint has distance at most \(5d\) from both
lifted balls. Any difference in this larger ball has length at
most \(10d<\ell/2\). Its Euclidean length therefore equals the
flat torus distance. This identifies the integral with the one
using the flat distance intrinsically. The lower separation follows
from \eqref{fp:eq-shape} and the upper separation from the diameter
of the common ball. Integration against the nonnegative product
gives the two middle inequalities. The mass floor and the volume
of each support ball give the outer inequalities. Under dilation
the distances are multiplied by \(t\), the Lipschitz constants by
\(t^{-1}\) and the masses by \(t^4\). The factor \(t^8\) in the
product measure cancels the homogeneity \(t^{-8}\) of the kernel.
\end{proof}

\subsection{Curvature normalization}

In the pairing of Section~\ref{fp:sec-setting}, take an orthonormal basis
of Hermitian traceless colour generators \(T_\alpha\),
\(1\le\alpha\le3\), with
\(\tfrac12\operatorname{tr}(T_\alpha T_\beta)=\delta_{\alpha\beta}\).
For a raw link increment \(A\), the leading plaquette logarithm
of \(e^{itA}\) is \(it\,\mathrm d A\), where \(\mathrm d\) is the
unscaled oriented link circulation. Its trace vanishes at every
order. Expanding the exponential in the Wilson action gives
\[
 1-\tfrac12\operatorname{Re}\operatorname{tr}U_p(t)
   =\tfrac12 t^2|(\mathrm dA)(p)|^2+O(t^3).
\]
Thus the free Wilson curvature at coupling \(g\) is the
three-colour Gaussian with covariance \(g^2P_{\rm ex}\) per
colour, where \(P_{\rm ex}\) is the counting-coordinate
orthogonal projection onto \(\operatorname{ran}\mathrm d\).
Indeed \(\mathrm d\) is an isomorphism on
\((\ker\mathrm d)^\perp\) onto this range, and its constant
Jacobian cancels in the normalized Gaussian integral.

For one colour, differentiating
\(\delta_{\mu\nu}/(4\pi^2|u|^2)\) once in each curvature argument
gives the unit-coupling continuum curvature matrix on
\(\Lambda^2\mathbb R^4\).
In detail, with \(G(u)=(4\pi^2|u|^2)^{-1}\) and
\(J=\widehat u\widehat u^{\mathsf T}\), its action is
the negative of \(\nabla^2G=(4J-I)/(2\pi^2|u|^4)\)
on each factor of the wedge, added together. Since \(J\)
has rank one, the term \(Jv\wedge Jw\) vanishes. Consequently
\begin{equation}\label{fp:eq-free-coefficient}
 \mathsf K(u)=\frac{\Lambda^2(I-2J)}{\pi^2|u|^4},\qquad
 \|\mathsf K(u)\|_{\mathrm{HS}}^2=\frac{6}{\pi^4|u|^8},\qquad
 C_G=\frac{9}{\pi^4}.
\end{equation}
The norm calculation uses orthogonality of \(I-2J\) and the six
positive plane labels. For real centred Gaussian coordinates \(X,Y\)
of any covariance, \(\operatorname{Cov}(X^2,Y^2)=2(\mathbb E XY)^2\).
The two half-density factors give \(1/4\), and summing over
independent colours gives \(3\). Thus the separated half-square
covariance at coupling \(g\) is \(C_Gg^4\mathcal K(f,h)\).
This calculation concerns only separated correlations.

\begin{proposition}[Free reference at a blocked scale]
\label{fp:prop-free-reference}
On the periodic fine lattice of spacing \(a=L^{-K}\) and side \(\ell\),
use a compatible block hierarchy of odd factor \(L\).
Take \(0\le j\le K\), \(b=L^ja\), and for \(j\ge1\) require at
least two sites per direction at mesh \(b\). Let \(P_j\) be
the unit-coupling counting-coordinate covariance of the
conditional mean of the exact free curvature given its
level-\(j\) plaquette sheet averages. In particular \(P_0=P_{\rm ex}\).
Use the barycentre \(x_p\) of each positively oriented fine
plaquette \(p\). Define the free reference sum by
\[
 \mathcal K_{a,j,\ell}(f,h)=
       \sum_{p,q}f(x_p)h(x_q)P_j(p,q)^2.
\]
Then the connected covariance of the centred free half-squares
of that conditional mean at coupling \(g>0\) is exactly
\(\tfrac32g^4\mathcal K_{a,j,\ell}(f,h)\).
For \((f,h)\in\mathfrak T_{T_\ell}(d;\vartheta)\) with
\(10a\le d\le\ell/100\) and \(b\le d\), it satisfies
\begin{equation}\label{fp:eq-free-reference-error}
 \left|\frac{(3/2)\mathcal K_{a,j,\ell}(f,h)}
                 {C_G\mathcal K(f,h)}-1\right|
 \le C_L\left[
     (a/d)^{1/2}+(d/\ell)^4+\frac{a}{\vartheta^3d}
     +\mathbf1_{\{j\ge1\}}e^{-c_Ld/b}\right].
\end{equation}
The constants are uniform in \(a,j,\ell,g\), positions and amplitudes.
\end{proposition}

\begin{proof}
The orthogonal splitting for the actual sheet averages in
\cite[Proposition~3.2]{GI} identifies \(P_j\). The Gaussian
square identity gives the exact covariance above. The counting
entries \(P_j(p,q)\) become physical entries after multiplication
by \(a^{-4}\), since the physical curvature is \(a^{-2}\) times
the raw flux. This change supplies \(a^8\) in the double
quadrature for the squared physical entries. In these conventions
\cite[Theorem~D]{GI} gives \eqref{fp:eq-free-reference-error}.
In its application here, every cross-displacement lies in
\([d,10d]\) by Lemma~\ref{fp:lem-shape-kernel}. Thus the
pointwise comparison range \([10a,\ell/10]\) contains every
contributing pair, including pairs whose support balls meet
a seam of the torus. The lost fluctuation kernel has the
exponential bound in that theorem at distance \(r\ge b\).
The shifted-grid quadrature has absolute cost
\(CaM_fM_h/(\vartheta d)\). Dividing by
\eqref{fp:eq-shape-kernel-bounds} gives the relative cost
\(Ca/(\vartheta^3d)\).
\end{proof}

\begin{lemma}[Bounds for the raw background and slice kernels]
\label{fp:lem-free-kernel-bounds}
Use the hierarchy and the counting-coordinate projections of
Proposition~\ref{fp:prop-free-reference}. At level \(j\) put
\(b_j=L^ja\), and let \(r\) be the flat torus distance between
the fine plaquette barycentres. Uniformly over the compatible
tori and \(0\le j\le K\),
\begin{equation}\label{fp:eq-free-point-bounds}
 \begin{split}
 |P_j(p,q)|&\le C_L(a/b_j)^4,\\
 |P_j(p,q)|&\le C_L a^4r^{-4}
       \qquad(\max(10a,b_j)\le r\le\ell/10).
 \end{split}
\end{equation}
For \(j\ge1\) the slice satisfies
\[
 |(P_{j-1}-P_j)(p,q)|
     \le C_L(a/b_{j-1})^4 e^{-c_L r/b_{j-1}} .
\]
\end{lemma}

\begin{proof}
Write \(n=L^j=b_j/a\). In raw counting coordinates let \(H_n B\)
be the least-norm fine exact curvature with sheet datum
\(\mathrm d_c B\), where \(\mathrm d_c\) is the coarse circulation.
If the fine torus has just one site, all raw circulations vanish,
so we may work with at least two sites at every used level.
The range of \(H_n\) is \(\operatorname{ran}P_j\) by the sheet identity and
the least-norm characterization. In fact the image of the sheet
map on exact fine curvatures is the range of \(\mathrm d_c\),
since the link average is onto. The lift bound proved in
\cite[proof of Proposition~6.4]{MB} gives
\[
 \sum_e |H_n(p,e)|\le C_L n^{-2}.
\]
Here \(e\) ranges over coarse bonds. The bound follows by
summing the kernel \(C_L n^{-2}e^{-c_L\operatorname{dist}(p,e)/b_j}\)
on the coarse lattice. The sum of its periodic images is
bounded by the sum on the full coarse lattice. The lift
estimate is at mesh \(b_j\) and uses no step above that mesh.
For \(n=1\) the lift is the ordinary plaquette circulation
and the row bound follows by counting its four entries.

Set \(\mathcal E_n=H_n^{\mathsf T}H_n\) and use \({}^+\) for the
inverse on \((\ker\mathrm d_c)^\perp\), extended by zero on
\(\ker\mathrm d_c\). The kernel is the same for \(\mathcal E_n\)
and \(\mathrm d_c\) by the defining datum and least-norm property.
By \cite[Proposition~4.1]{MB},
\(\mathcal E_n\ge\mathrm d_c^{\mathsf T}\mathrm d_c\).
Order reversal on the quotient gives
\(\mathcal E_n^+\le(\mathrm d_c^{\mathsf T}\mathrm d_c)^+\).
If there are \(N=\ell/b_j\ge2\) sites per direction, the
diagonal entries of the latter are bounded by
\[
 \frac1{N^4}\sum_{k\ne0}
       \frac1{4\sum_{\mu=1}^4\sin^2(\pi k_\mu/N)}
 \le \frac C{N^2}\sum_{0<|k|_\infty\le N/2}|k|^{-2}
 \le C .
\]
The frequencies are represented with \(|k_\mu|\le N/2\).
The eigenvalue displayed is the positive curl eigenvalue on
the Fourier complement of the discrete gradients, and the
zero frequency is removed. The last estimate counts at most
\(C m^3\) frequencies with \(|k|_\infty=m\).
Positivity and Cauchy--Schwarz then bound every entry of
\(\mathcal E_n^+\) by \(C\). Apply the row bound twice to
\(P_j=H_n\mathcal E_n^+H_n^{\mathsf T}\). This proves the first
line of \eqref{fp:eq-free-point-bounds}. All matrices here
use raw flux and counting pairings.

The slice bound is \cite[Proposition~6.4]{MB}, expressed in the
distance \(r\) by moving the plaquette base points by \(O(a)\).
At the distances of the second line, \cite[Proposition~8.3]{GI}
and \eqref{fp:eq-free-coefficient} give \(|P_0(p,q)|\le Ca^4r^{-4}\).
For \(1\le i\le j\), put \(t_i=r/b_{i-1}\).
The sum of the slice bounds is at most
\(C_L a^4r^{-4}\sum_i t_i^4e^{-c_Lt_i}\).
Since \(r\ge b_j\), all \(t_i\ge1\). Use
\(t_i^4e^{-c_Lt_i}\le C_L/t_i\) and the geometric sum of
\(b_{i-1}/r\) to bound the last sum uniformly. Subtracting the
slices from \(P_0\) proves the second line.
\end{proof}


\section{The witness scale and the source coefficient}
\label{fp:sec-witness}

Fix a terminal floor \(H>1\) and a first-passage trajectory
\(x_0=\beta,\ldots,x_K\) of the real construction in
\cite[Theorem~A]{RT}. Use the same formed coefficients in every
admitted finite realization of that prefix. Thus
\(H<x_K\le H+B\), and there are fixed \(\lambda,B>0\) and \(C_0\ge0\)
such that
\begin{equation}\label{fp:eq-depth-drift}
 \lambda(j-i)-C_0\le x_i-x_j\le B(j-i)
 \qquad (0\le i\le j\le K).
\end{equation}
The actual increments have a common one-loop part \(b_i^0\),
whose partial sums differ from \(ic_0\) by at most a fixed
constant, and a remaining part bounded at step \(i\) by
\(e(x_{i-1})\), where the nonnegative envelope \(e\) is bounded on \([H,\infty)\)
and tends to zero at infinity.
In the half-trace Wilson convention,
\[
 c_0=\frac{11}{3\pi^2}\log L.
\]
These are the real trajectory and coefficient inputs of
\cite{RT}. The one-loop coefficient is a determinant and trace
input in that construction.

Write \(\eta_j=L^{j-K}\) for the physical mesh at level \(j\), and
let \(M\) be the fixed side of a localization cube in its own
mesh units. Fix the shape parameter \(\vartheta>0\) and the
positive logarithmic exponent \(p\) of the small-field radius.
Define
\begin{equation}\label{fp:eq-witness-room}
 \overline x_s=H+B(s+1),\qquad
 D(s)=D_0(\log\overline x_s)^{2p},\qquad
 s(d)=\min\{s\in\mathbb Z_{\ge0}:L^sd\ge D(s)\}.
\end{equation}
The fixed constant \(D_0\) is chosen after \(L,M,\vartheta\), with
room for the localization and source ranges specified below.
Choose it so that \(D(s)\ge1\). This can be done uniformly for
\(H\ge e\).

\fpneedspace{25\baselineskip}
\begin{proposition}[Choice of scale]\label{fp:prop-witness}
For \(0<d<1\), let \(s=s(d)\) and use a cutoff with \(K\ge s\).
Set
\[
 k_0=K-s,\qquad a_s=L^{-s},\qquad
 \widehat d=d/a_s,\qquad
 x=x_{k_0},\qquad g_K(d)=x^{-1/2}.
\]
Then \(D(s)\le\widehat d<LD(s)\), and at fixed \(H\),
\begin{equation}\label{fp:eq-witness-asymptotic}
 \frac{s(d)\log L}{\log(1/d)}\longrightarrow1,\qquad
 \frac{g_K(d)^{-2}}{\log(1/d)}
   \longrightarrow c_*:=\frac{11}{3\pi^2}
 \qquad(d\downarrow0),
\end{equation}
uniformly over the admitted first-passage data with \(K\ge s(d)\).
In particular \(\widehat d\) is bounded above and below by
fixed positive multiples of \((\log x)^{2p}\) for small \(d\).
If
\[
 m=\lceil 2\log_L\overline x_s\rceil,\qquad k_1=k_0+m,
\]
then \(k_1\le K\) for small \(d\), and
\(\sup_{k_0\le j\le k_1}|x_j/x-1|\to0\) uniformly in the same data.

Suppose the complete clean reads at these levels have a joint lift of
diameter at most \(P(s)L^{-(s-m)}\), where \(P\) is a fixed
polynomial with nonnegative coefficients. On any fixed reference
torus they fit in a common injective patch for all small \(d\).
For a shape test \(f\), its ordinary source norm at level \(j\) is
\[
 \|f\|_{\sharp,K-j}=\max(\|f\|_\infty,c_M\eta_j\operatorname{Lip} f),
\]
where \(c_M\) is the fixed locality constant of \cite{SI}.
For \(D_0\) sufficiently large this is \(\|f\|_\infty\) at \(j=k_0\),
and is at most \(L^{j-k_0}\|f\|_\infty\) for \(k_0\le j\le k_1\).
\end{proposition}
\begin{proof}
Since \(\log D(s)=O_H(\log\log(s+2))\), the exponential
\(L^s d\) eventually exceeds \(D(s)\). The minimum is positive
when \(d<1\). Its defining inequality and the failure at \(s-1\)
give the two bounds on \(\widehat d\), using that \(D\) is
nondecreasing. They give
\(s\log L=\log(1/d)+O_H(\log\log(s+2)+1)\).
Also \(s\to\infty\) as \(d\downarrow0\), since each fixed set of
candidate levels fails for small enough \(d\). The first limit
follows on dividing by \(s\log L\).

For the second limit, write \(e_*(y)=\sup_{t\ge y}e(t)\) on the
admitted half-line. It is bounded by hypothesis.
By \eqref{fp:eq-depth-drift},
\(x_{K-u}\ge H+\max(0,\lambda u-C_0)\).
Telescoping the last \(s\) increments yields
\[
 |x_{K-s}-x_K-c_0s|
 \le C+\sum_{u=1}^s e_*\bigl(H+\max(0,\lambda u-C_0)\bigr)=o(s).
\]
The last bound follows by averaging a bounded sequence tending to
zero, independent of the datum. The bounded terminal interval
and the first limit give \eqref{fp:eq-witness-asymptotic}.
The lower drift and the upper envelope \(\overline x_s\) give
\(x\asymp_H 1+s\) and \(\overline x_s\asymp_H x\).
This gives the stated size of \(\widehat d\).

At fixed \(H\), \(m=O_H(\log(s+2))\) and \(L^m\le L\overline x_s^2\).
The drift bounds give \(|x_j-x|\le Bm+C_0=o(x)\) through \(k_1\).
The lifted diameter tends to zero since
\(P(s)L^{m-s}\le L P(s)\overline x_s^2L^{-s}\).
When it is less than a quarter of the reference period, a ball
about a point of the read contains the lift and projects
injectively. This is a geometric consequence of the supplied
bound on the complete reads. Finally, the shape condition gives
\(c_M\eta_{k_0}\operatorname{Lip} f
 \le c_M\|f\|_\infty/(\vartheta\widehat d)\).
Choose \(D_0\) so the last prefactor is at most one.
The bound at \(j\ge k_0\) then follows from
\(\eta_j=L^{j-k_0}\eta_{k_0}\).
\end{proof}

\begin{corollary}[One-loop comparison and dilation]
\label{fp:cor-witness-variants}
At fixed \(H\), let
\[
 g_{\rm lin}(d)=\bigl(H+c_*\log(1/d)\bigr)^{-1/2}.
\]
Uniformly on the first-passage data of \cref{fp:prop-witness},
\(g_K(d)/g_{\rm lin}(d)\to1\).
In a physical comparison with main factor \(g_K(d)^4\), this
replacement multiplies the existing relative factor by
\((g_K(d)/g_{\rm lin}(d))^4\).

Let \(K'\) be the first-passage depth for the same bare coupling
\(\beta\), and let \(b=K-K'\). The Wilson measure on the torus
of side \(\ell\) at cutoff \(K\) is carried by lattice dilation
to the torus of side \(L^b\ell\) at cutoff \(K'\).
For each test put \(f^{[b]}(x)=f(L^{-b}x)\). Then
\[
 S^T_{2;K,\ell,\beta}(f,h)=
 S^T_{2;K',L^b\ell,\beta}(f^{[b]},h^{[b]}).
\]
The shape parameter and the free integral \(\mathcal K\) are
preserved, with observation scale \(L^b d\) on the right.
A first-passage comparison therefore applies on the right
when the dilated datum satisfies its volume, scale and cutoff
hypotheses.
\end{corollary}
\begin{proof}
By \eqref{fp:eq-witness-asymptotic}, the two inverse squared
couplings divided by \(\log(1/d)\) have the same positive limit.
This gives the ratio and the stated relative factor.

Dilation by \(L^b\) maps spacing \(L^{-K}\) to \(L^{-K'}\) and
bijects oriented links and plaquettes. It preserves each
plaquette product and each Haar factor. The normalized
measures are therefore identified at the same \(\beta\).
The transformed test on each corresponding barycentre gives
the same observable. The support radii and distance multiply
by \(L^b\), each mass by \(L^{4b}\), and each Lipschitz constant
by \(L^{-b}\). In the free integral the two volume factors
cancel the scaling of \(|x-y|^{-8}\). These equalities hold
also for negative \(b\).
\end{proof}

\subsection{A constant source in the fixed reference}

Freeze the reference data of the unsourced trajectory at \(\beta\),
including \(\theta_i=x_i\), the cubes, the cuts, and the readout
maps. On a constant interior read, let \(b_i(v)\) be the actual
Wilson coefficient of the complete formed member at step \(i\)
when the initial scalar source is \(v\). The coefficient is read
at its own pointing before its Wilson subtraction. The formation
includes the marked contribution from earlier retained remainders.
The coefficients are those of the common companion, so the
interior scalar values agree on the admitted volumes.
In particular \(b_i(0)=x_{i-1}-x_i\).

We use the fixed-reference analytic and real matching hypotheses of
\cite[Propositions~4.1 and~4.2]{SI}. For the analytic part, the
ordinary coefficient is the same formed holomorphic map of the
old weight deviations in the sourced and unsourced presentations.
More precisely, it has the form
\[
 b_i(v)=b_i^0+\mathcal B_i\bigl(
        (\theta_l-x_l+Z_l(v))_{0\le l<i}\bigr)+R_i(v),
\]
where \(b_i^0\) is source independent, \(Z_l\) is the shift
defined in \eqref{fp:eq-scalar-shift}, and \(R_i\) is the marked
part. The ordinary map \(\mathcal B_i\) has no additional
dependence on \(v\). It is holomorphic on the full relative
product of weight discs. The half product has room
\(6\varrho\) beyond the real base weights, for a fixed source radius
\(\varrho>0\). On the half product its absolute derivative row sum is
bounded by
\(C x_{i-1}^{-3/2}(\log x_{i-1})^p\).
This is the regular coefficient bound on the actual creation
domains. The marked coefficient is holomorphic and bounded by
\(C x_{i-1}^{-p_R}\) on \(|v|<4\varrho\), with \(p_R>1\). The
coefficient maps preserve reality. For matching, the real weights have \(2\varrho\) of room in the
common tube and the full increments beyond the common one-loop
part are bounded by \(\varrho/3\). The shifted bare input is
admitted and each subsequent weight in that tube has actual local
successor admission. The two real
presentations have common hereditary reads, centered formation,
real cuts, vacuum convention and cofinal physical coefficient
readout as in the same-frame source theorem of \cite{SI}.
The numerical room does not supply these local admission tests.

\begin{proposition}[The physical source factor]
\label{fp:prop-source-normalization}
Under these hypotheses, above a fixed floor for \(H\), set
\begin{equation}\label{fp:eq-scalar-shift}
 Z_j(v)=v+\sum_{i=1}^j\bigl(b_i(v)-b_i(0)\bigr),\qquad
 \zeta_j=Z_j'(0).
\end{equation}
The constant-source weight at level \(j\) is \(x_j-Z_j(v)\).
For the shifted real input, through the same fixed prefix,
\[
 x_j^{\rm fix}(\beta-v)=x_j-Z_j(v) \qquad (|v|<\varrho).
\]
In particular
\begin{equation}\label{fp:eq-source-jacobian}
 \zeta_j=\left.\frac{d}{db}x_j^{\rm fix}(b)\right|_{b=\beta}
        =1+\sum_{i=1}^j b_i'(0),\qquad
 \sum_{i=1}^j |b_i'(0)|
 \le C\bigl[x_j^{-1/2}(\log x_j)^p
                   +\varrho^{-1}x_j^{1-p_R}\bigr].
\end{equation}
The floor can be chosen so that \(2/3\le\zeta_j\le4/3\).
At the witness set \(\zeta_K(d)=\zeta_{k_0}\). Then
\(\zeta_K(d)\to1\) as \(d\downarrow0\), uniformly over the
admitted first-passage data with \(K\ge s(d)\).
If \(p_R\ge3/2\), its difference from one is bounded by
\(C g_K(d)(\log g_K(d)^{-2})^p\).
\end{proposition}
\begin{proof}
At the bare level the source changes the coefficient in the Wilson
weight from \(\beta\) to \(\beta-v\). The coefficient subtracted
from the regular member is added to the exponent of that weight
as its own pointed Wilson action. Thus the inverse weight loses
\(b_i(v)\) at step \(i\). The stage sums of tents equal one on
the interior read. The resulting weight is
\(\beta-v-\sum_{i\le j}b_i(v)=x_j-Z_j(v)\).
For real \(v\), apply the same-frame source theorem of \cite{SI}
with the stated hereditary admission and common readout. This
identifies that weight with the real shifted presentation while
the reference data stay fixed. The derivative in
\eqref{fp:eq-source-jacobian} follows by the chain rule.

For the estimate put
\(d_i=\sup_{0<|v|<4\varrho}|b_i(v)-b_i(0)|/|v|\).
Assume inductively that \(\sum_{l<i}d_l\le1/6\). Each previous
weight shift then has modulus at most \((7/6)|v|<6\varrho\).
The segment between the unsourced and sourced old weights lies
in the half product. Integrating its derivative row bound
bounds the unmarked difference quotient by
\(C x_{i-1}^{-3/2}(\log x_{i-1})^p\).
Subtracting the value at zero from the marked coefficient and
applying Schwarz's lemma on \(|v|<4\varrho\) bounds its difference
quotient by \(C\varrho^{-1}x_{i-1}^{-p_R}\). The one-loop part is
common and cancels.

For \(i\le j\), the lower drift gives
\(x_{i-1}\ge x_j/2+\lambda(j-i+1)\) once \(H\ge2C_0\).
Both envelope functions decrease above a fixed floor.
For \(A\ge e\) and \(q>1\), their required integral bound is
\[
 \int_A^\infty y^{-q}(\log y)^p\,dy
 =A^{1-q}\int_0^\infty e^{-(q-1)t}(\log A+t)^p\,dt
 \le C_{p,q} A^{1-q}(\log A)^p .
\]
Comparison of the sums with this integral gives the right side
of \eqref{fp:eq-source-jacobian} as a bound on every radial sum.
Increasing \(H\) makes this bound at most \(1/6\) throughout the
prefix. This closes the induction, starting with \(Z_0(v)=v\).
The derivative at zero satisfies \(|b_i'(0)|\le d_i\).
Reality yields the claimed interval for \(\zeta_j\). The witness
limit follows from \(x\to\infty\) in
\eqref{fp:eq-witness-asymptotic}, and the last bound follows from
\(x^{1-p_R}\le x^{-1/2}\) when \(p_R\ge3/2\).
\end{proof}

\subsection{Differentiating the pointed classical part}

On an ordinary interior read at the witness, all the used stage
profiles equal one. Here the action is evaluated on the background
formed from the level-\(k_0\) data, using its bare-mesh plaquettes.
The witness mesh \(a_s\) controls the read range. For each step \(j\le k_0\), let \(h_y^{(j)}(p)\)
be its nonnegative tent weights at the plaquette \(p\) of the
classical read and let \(p_y\) be its fixed bare source anchor.
The exact classical identification in \cite{SI} transports the
pointed action to this read, with
\(\sum_y h_y^{(j)}(p)=1\) and
\(\operatorname{dist}(x(p),x(p_y))\le C M\eta_j\)
when \(h_y^{(j)}(p)\ne0\). These statements apply on the complete
common read, with its projected branch tests and contours.
They concern the literal transported Wilson counterterms.

\begin{proposition}[Pointed first derivatives]
\label{fp:prop-pointed-insertion}
For \(W=z_1f_1+z_2f_2\) with disjointly supported shape tests,
define \(f_i^w\) on the bare-mesh plaquettes of the witness background by
\begin{equation}\label{fp:eq-physical-profile}
 \zeta_{k_0} f_i^w(p)=f_i(x(p))+
 \sum_{j=1}^{k_0}b_j'(0)\sum_y h_y^{(j)}(p) f_i(x(p_y)).
\end{equation}
The first derivative of the sourced classical exponent at zero
is the Wilson action weighted by \(\zeta_{k_0}f_i^w\). The mixed
derivative of that classical exponent vanishes. The discrepancy
and support obey
\begin{equation}\label{fp:eq-profile-range}
 \begin{split}
 |f_i^w(p)-f_i(x(p))|
 &\le C M a_s\operatorname{Lip}(f_i)
                 \sum_{j\le k_0}|b_j'(0)|,\\
 \{x(p):f_i^w(p)\ne0\}
 &\subset (\operatorname{supp}f_i)^{+CMa_s}.
 \end{split}
\end{equation}
Here \(E^{+r}\) denotes the closed \(r\)-neighbourhood of \(E\).

For a bilinear form on these plaquette weights, with the physical
quadrature at the bare mesh \(a=L^{-K}\), write
\[
 \mathcal B(u,v)=a^8\sum_{p,q} u(p)v(q)\,B(p,q),
\]
assume \(|B(p,q)|\le A d^{-8}\) between the two enlarged supports.
When \(C M a_s\le d/4\), replacing both sampled tests by the
profiles \(f_i^w\) changes this form by at most
\begin{equation}\label{fp:eq-profile-kernel-error}
 C A\frac{M}{\vartheta\widehat d}
        \left(\sum_{j\le k_0}|b_j'(0)|\right)
        \|f_1\|_\infty\|f_2\|_\infty.
\end{equation}
The profile can have either sign.
\end{proposition}
\begin{proof}
At the plaquette \(p\) the source changes the inverse weight by
minus
\[
 W(x(p))+\sum_{j\le k_0}\sum_y h_y^{(j)}(p)
             \bigl[b_j(W(x(p_y)))-b_j(0)\bigr].
\]
The weights and anchor locations are fixed. This gives
\eqref{fp:eq-physical-profile}. Each summand depending on a
single anchor has mixed derivative proportional to
\(f_1(x(p_y))f_2(x(p_y))=0\). The bare source is linear.

For the exact accounting, let \(G_{j,y}(v)\) be the unnormalized
pointed regular member on the fixed read and \(U\) its classical
field. Its normalized exponent contains
\[
 G_{j,y}(v)-G_{j,y}(v)(\mathbf1,0)
                    -b_j(v)A(h_y^{(j)},U),
\]
where \(A(h,U)\) denotes the Wilson action with plaquette weight
\(h\) and the exact inherited transport. The difference of the
last term between \(v=W(x(p_y))\) and \(v=0\) cancels precisely
the contribution of the same coefficient to the shifted weight
on the common read. Thus differentiating the bare source and
the regular members before their Wilson subtractions gives the
same answer as differentiating the running weight and the
already subtracted members. In the latter representation,
\eqref{fp:eq-physical-profile} is the classical contribution.
Its paired regular derivative is the derivative of the
normalized member, with its own Wilson subtraction. Missing
shares on other retained reads are separate terms in the
complete exponent.

Subtract \(\zeta_{k_0} f_i(x(p))\) on the right side of
\eqref{fp:eq-physical-profile} and use the unit sum of the tents.
The difference at each anchor is bounded by
\(CM\eta_j\operatorname{Lip}(f_i)\). Since \(\eta_j\le a_s\)
and \(\zeta_{k_0}\ge2/3\), this proves the first inequality in
\eqref{fp:eq-profile-range}. The range of the anchor read proves
the second one.

The widened supports have separation at least \(d/2\). Each
contains at most \(C(d/a)^4\) plaquette barycentres, with
multiplicity over the six orientations. This bound follows
by covering the enlarged ball by mesh cubes, with a bounded
number of barycentres per cube. The first inequality in
\eqref{fp:eq-profile-range}, the shape Lipschitz bound and the
bound on the slope sum imply \(\|f_i^w\|_\infty\le C_\vartheta
\|f_i\|_\infty\) there. Expand the difference of products into
\((f_1^w-f_1)f_2^w+f_1(f_2^w-f_2)\), with the original tests
evaluated at the respective barycentres. In the estimate on
each sum the factor \(a^8d^{-8}\) pays for both support
counts. This yields \eqref{fp:eq-profile-kernel-error}.
\end{proof}


\subsection{Summing the ordinary first source derivatives}
\label{fp:subsec-ordinary-response}

In this subsection set \(n=k_0\), \(x=x_n\), and let \(f\) be one of
the shape tests. Write \(M_f=\|f\|_\infty\) and \(M_i=\|f_i\|_\infty\)
for their amplitudes. The physical meshes at levels \(j\le n\) are
\(\eta_j=L^{j-K}\). Use the radii and the small-field profile of the
actual source formation in \cite{SI}, with
\[
 \delta_j=x_j^{-1/2}(\log x_j)^p,\qquad
 \alpha_j=x_j^{-1/2}(\log x_j)^q,\qquad s_j=\delta_j/\alpha_j.
\]
Fixed radius constants can be absorbed into the estimates below.
Take \(q>21p\) and \(p>1\). Formed regular and remainder
origins have indices \(1\le j\le n\); their lists are empty
when \(n=0\). On a support neighbourhood, assign each ordinary
source term once
in its own class and at its own root. A short ordinary read uses
the same localized field and current as its whole-lattice
formation at that root. The complete determining set includes
every branch test, chart read, gauge contour and source anchor.
This common-read condition is required before replacing a
history member by its formed jet. A local curvature test alone
does not identify those two evaluations.

\fpneedspace{37\baselineskip}
\begin{lemma}[The summed ordinary source bounds]
\label{fp:lem-ordinary-response-sums}
Suppose the localized backgrounds at the witness are real and
the passed test gives
\[
 |\log U_p|\le C(a/a_s)^2\delta_n
\]
on the bare plaquettes of the localized read. Suppose the
ordinary source formation and the hypotheses of
\cite[Proposition~5.6]{SI} hold on all the following reads.
\begin{enumerate}
\item The near regular and defect terms are filed in witness
windows meeting a fixed \(CMa_s\) enlargement of the support.
In each window their evaluated complete field and current datum
has size at most \(C\delta_n\) and satisfies the affine comparison
and truncation hypotheses of \cite[Proposition~5.6]{SI}.
At the older origins
\(\|f\|_{\sharp,K-j}\le M_f\) by the witness room.
\item For an ordinary remainder of origin \(j\), its actual
source family is holomorphic on its radius
\(\varrho/\|f\|_{\sharp,K-j}\). Its bound there is
\(C x_j^{-\kappa_0/2}e^{-\kappa l}\), where \(l\) is the length
of its domain and \(\kappa_0>4\). This is a bound both on its
admitted evaluated pairs and on the product of that source disc
with an old pair ball of radius \(c\alpha_j\) about the trivial
pair, on which the family is jointly holomorphic. The source
family is invariant under constant colour rotations in that chart.
\item A homogeneous nonwrapping read of age \(A=n-j\) and length
\(l\) has old pair norm at most
\(C(1+l)^b\delta_n L^{-2A}\) in its admitted axial frame.
Its current is included in the pair norm. The ordinary
long regular jets on these reads have the same admitted core and
trivial-pair ball bounds as their weighted first-source amplitude.
They vanish with their first field derivative at the trivial pair.
Those outside the near windows have \(l\ge cL^A-C\).
\item At origin \(j\) the domain count at one root and length \(l\)
is at most \(C e^{\kappa_D l}\), with \(\kappa_D<\kappa\). The complete
determining set together with its root has diameter at most
\(C(1+l)\eta_j\).
For a nonhomogeneous ordinary read which reaches a distant live
region or winds on the reference torus, require
\(l\ge c\widehat d L^A-C\). Use the unsubstituted source-disc bounds
on those reads, with a fixed polynomial in \(x,L^A\) allowed
in their individual majorants at rate \(\kappa\).
The ordinary two-source terms are holomorphic with these
absolute majorants on a fixed source bidisc of radii proportional
to \(M_i^{-1}\) and
depend only on the source values in their declared complete
determining sets. They include the terms in the class-independent
scalar source energy before vacuum subtraction, as well as the
remaining normalized interactions. The local scalar energy uses
the same complete source reads and bounds, as in
\cite[Proposition~7.1]{SI}.
\end{enumerate}
All constants in these inputs are uniform in the cutoff and
actual retained histories. Let \(I(f)\) be the sum of the
ordinary normalized first-source terms just specified, and
\(X(f)\) the localized Wilson read with profile from
\eqref{fp:eq-physical-profile} and its passed-test indicator.
For any fixed \(b_0\) with \(0<b_0<\kappa_0/2-2\), the bounds are
\[
 \begin{split}
 \|X(f)\|_\infty
   &\le Cx^{-1}\widehat d^{\,4}(\log x)^{2p}M_f,\\
 \|I(f)\|_\infty
   &\le Cx^{-1}\widehat d^{\,4}
             [(\log x)^{3p-q}+x^{-b_0}]M_f .
 \end{split}
\]
The ordinary mixed terms are bounded by
\(C\widehat d^{\,4} e^{-c'\widehat d}M_1M_2\).
Terms filed as live, missing shares of the source number, and
errors replacing the history background by its localization
use separate bounds.
\end{lemma}

\begin{proof}
At a passed-test plaquette the half-trace action is bounded
by \(C(a/a_s)^4\delta_n^2\). The profile and support estimate
\eqref{fp:eq-profile-range} bounds its weight by \(CM_f\)
and the number of its bare plaquettes by \(C(d/a)^4\).
Multiplication gives the bound for \(X(f)\).

The near-window bound in \cite[Proposition~5.6]{SI} gives
\(C\delta_n^2s_n M_f\) per witness window. It is used with the
same source anchor for its complete pointed member and its
own Wilson subtraction. At most \(C\widehat d^4\) such windows
meet the enlarged support. The contribution of this part is
therefore bounded by \(C\widehat d^4 M_f x^{-1}
(\log x)^{3p-q}\).

For the long regular terms, \cite[Proposition~5.6]{SI} gives
\(A_j\le C\alpha_{j-1}^2s_{j-1}\). Its pointed coefficient bound
keeps the same norm after the member's own Wilson subtraction.
The adjacent drift bounds give \(\alpha_{j-1}\le C\alpha_j\), since
\(|x_{j-1}-x_j|\le B+C_0\). The lower drift gives
\(x_{j-1}\ge x_n-C_0\), so \(s_{j-1}\le C s_n\) since \(q>p\).
Using the older sharp norm bound in the statement, the regular
jet majorant at these reads is thus
\(C\alpha_j^2s_n M_f e^{-\kappa l}\).

For a homogeneous old domain, use its field and current ratio
\[
 u_j=C(1+l)^b(\delta_n/\alpha_j)L^{-2A}.
\]
For an old remainder, source Cauchy and vacuum subtraction
give the jet majorant
\[
 C\varrho^{-1}M_f x_j^{-\kappa_0/2}e^{-\kappa l}
\]
on the old pair ball and on the evaluated pair. The subtraction
kills the constant term. Its linear field term vanishes because the
adjoint representation has no invariant linear functional.
On the half ball of its old pair chart, Cauchy's formula
bounds the evaluated jet by \(C u_j^2\) times its jet
majorant. When \(u_j\) is greater than a fixed small threshold,
the bound on the actual admitted core gives the same estimate
after increasing the constant. This also applies to each
homogeneous long regular jet, whose two vanishing terms are
part of its normalized formation. No small field ratio is
asserted on nonhomogeneous reads.

The radius-\(d\) support and the determining-set bound give
at most \(C(\widehat d L^A+1+l)^4\) possible roots at length
\(l\). Indeed a root of a read meeting the support is within
distance \(d+C(1+l)\eta_j\) of its centre, and
\(d/\eta_j=\widehat d L^A\). Use the bound
\(C\widehat d^4L^{4A}(1+l)^4\).
Summing \(u_j^2\) times the remainder majorant over lengths
with the strict entropy margin cancels \(L^{-4A}\) against
\(L^{4A}\). The resulting sum is at most
\[
 C\widehat d^4 M_f\delta_n^2
            \sum_{j\le n} x_j^{-\kappa_0/2}\alpha_j^{-2}.
\]
Since \(\alpha_j^{-2}\le Cx_j\), comparison with an integral
using \(x_j\ge x/2+\lambda(n-j)\) bounds the sum by
\(C x^{2-\kappa_0/2}\). Its product with \((\log x)^{2p}\)
is \(O(x^{-b_0})\). The homogeneous long regular jets use
\(\alpha_j^2s_n\) in place of the remainder factor.
Their lower length bound leaves \(e^{-c'L^A}\) after summing
the length polynomial at the spare rate. Summing this in
\(A\) gives \(C\widehat d^4M_f\delta_n^2s_n\).

For a nonhomogeneous or winding read, sum the original
majorant over its length lower bound. The resulting term
has an \(e^{-c'\widehat d L^A}\) factor. The root count and
any fixed polynomial in \(L^A,x,\widehat d\) are absorbed
by a portion of this exponential, uniformly in age, because
\(\widehat d\) grows as \((\log x)^{2p}\).
An ordinary term is independent of either source variable whose
support misses its read. Thus the contributing mixed terms reach
both supports and their lengths have the same lower bound.
Cauchy's formula on the two source circles gives \(C M_1M_2\)
times the individual analytic bound. The same root and age
summation proves the mixed estimate. The classical mixed
terms vanish by \cref{fp:prop-pointed-insertion}.
\end{proof}


\section{The covariance on the reference torus}
\label{fp:sec-reference-comparison}

All the source derivatives in this section use the same bare datum and
the same fixed auxiliary choices as in the unsourced run. We write
$j_0=k_0$, $j_1=k_1$, $n=j_1-j_0$, $x=x_{j_0}$, $g_x=x^{-1/2}$,
and $\widehat d=d/a_s$ for the witness variables of
\cref{fp:prop-witness}. Here $a_s=L^{-s}$ is the witness mesh and
$a=L^{-K}$ is the bare mesh. Put $M_i=\|f_i\|_\infty$ for the two
source profiles. A positive law on a retained state includes the
standing class, the block field and all still stored slots. We keep
these slots when differentiating source integrals.

\begin{lemma}[Derivatives of the actual source integral]
\label{fp:lem-exact-insertion}
At a stop \(j\), suppose the actual summed source integral is
\[
 \frac{Z(z_1f_1+z_2f_2)}{Z(0)}=\mathbb E_j \mathscr W_j(z),
 \qquad \mathscr W_j(0)=1\quad\hbox{almost surely},
\]
where \(\mathbb E_j\) uses the unsourced positive retained-state
law. Assume \(\mathscr W_j\) is measurable and pointwise holomorphic on
a common source bidisc, and integrably dominated on every
compact subdisc. Define its logarithmic jets at zero by
\[
 \Phi_i=\partial_i\mathscr W_j(0),\qquad
 \Phi_{12}=\partial_1\partial_2\mathscr W_j(0)-\Phi_1\Phi_2.
\]
Assume the products used below are integrable. Then the bare
connected covariance satisfies
\begin{equation}\label{fp:eq-exact-covariance}
 S_2^T(f_1,f_2)=\operatorname{Cov}_j(\Phi_1,\Phi_2)
                         +\mathbb E_j\Phi_{12}.
\end{equation}
If the actual retained weight is locally \(e^\Phi\), these are
its first and mixed exponent derivatives. No common zero-free
disc for all retained weights is required.
\end{lemma}
\begin{proof}
The compact-disc majorants and Cauchy's formula permit
derivatives of \(\mathscr W_j\) to pass under the integral. Its integral
equals one at zero. The mixed derivative of its local logarithm
is therefore
\[
 \mathbb E_j\partial_{12}\mathscr W_j(0)
       -\mathbb E_j\Phi_1\,\mathbb E_j\Phi_2.
\]
Substitute the definition of \(\Phi_{12}\). The bare integral
gives the physical covariance by differentiating its bounded
insertions. Each individual retained weight is nonzero on some
neighbourhood of zero, so its local logarithmic derivatives
agree with the displayed jets. That neighbourhood need not be
uniform in the retained state. Integrating stored variables
before computing these jets adds their conditional covariance
to the mixed logarithmic derivative. The use of their actual
retained source weight preserves this contribution.
\end{proof}

At the witness, write \(\zeta=\zeta_K(d)\) and
\(\widetilde f_i=f_i^w\) for the profiles in
\eqref{fp:eq-physical-profile}. Proposition~\ref{fp:prop-pointed-insertion}
gives their once-counted classical derivative and the vanishing
of its mixed single-anchor terms.

\begin{lemma}[Removal of the other witness source derivatives]
\label{fp:lem-witness-errors}
Use \cref{fp:lem-exact-insertion} at $j_0$. Suppose the actual source
filing gives
\[
 \Phi_i=c_i+\zeta X_i+I_i+B_i ,
\]
where $c_i$ is the same scalar on all standing classes and
$X_i$ is the background Wilson read with weight $\widetilde f_i$
and its fixed local passed-test indicator. The scalar is the
derivative of the complete scalar source energy. Its mixed
derivative is included among the ordinary two-source terms of
\cref{fp:lem-ordinary-response-sums}. Missing shares on a retained
class are included in $B_i$. Assume the
complete ordinary read, response and summation hypotheses of
\cref{fp:lem-ordinary-response-sums}, with $p>1$ and $q>21p$.
Suppose $\widehat d\le C(\log x)^{2p}$ and $|\zeta|\le4/3$.
Assume $B_i$ and the nonordinary mixed derivative consist of
tails bounded, after division by the same profile norms, by fixed
powers of $\log x$ times
$e^{-c\widehat d}+e^{-c\check R}$, and of terms supported on
contact with a live component.
On that event their respective pointwise majorants, after division
by $M_i$ and $M_1M_2$, are
\[
 C (\log x)^C L^{C\check R}(1+A)^C L^{4A},
 \qquad
 C (\log x)^C L^{C\check R}(1+A)^C L^{8A},
\]
where $A$ is the largest contacting lineage age. Suppose there are
at most $C\widehat d^4$ possible contact roots and the actual
age tails at each root are bounded by
$C e^{-c(\log x)^p}(2L^8)^{-a}$ for $A\ge a$.
Assume $\check R/\log x\to\infty$,
$\check R=o((\log x)^p)$ and
$\widehat d\ge c(\log x)^{2p}$.
Then, uniformly for these data,
\begin{equation}\label{fp:eq-witness-errors}
 S_2^T(f_1,f_2)=\zeta^2\operatorname{Cov}_{j_0}(X_1,X_2)
                          +o(x^{-2})M_1M_2.
\end{equation}
\end{lemma}

\begin{proof}
The ordinary mixed derivative has no pointwise classical part.
The remaining ordinary mixed derivative is bounded by
$C\widehat d^4 e^{-c\widehat d}M_1M_2$ by
\cref{fp:lem-ordinary-response-sums}. The room makes it
smaller than any inverse power of $x$. The same lemma gives
the sup bounds for $X_i$ and $I_i$, with some $b>0$ in its
inverse-power error.

On the contact event the age estimate gives
\[
 \mathbb E_{j_0}\big[1_{\mathrm{contact}}(1+A)^C L^{8A}\big]
 \le C\widehat d^4 e^{-c(\log x)^p}
                  \sum_{a\ge0}(1+a)^C2^{-a}.
\]
We used the union bound over possible roots and the age tail at
the same lineage whose fields appear in the pointwise majorant.
This bounds $\|B_i\|_2/M_i$ and the mean nonordinary mixed
derivative by $o(x^{-N})$ for every fixed $N$. The loss
$L^{C\check R}$ is smaller than the rarity exponential, and
both far exponentials are superpolynomial by the lower
room and localization-radius bounds.
Cauchy--Schwarz controls all covariance terms with a $B_i$.
The constants $c_i$ have zero covariance with every variable.
Using $|\operatorname{Cov}(Y,Z)|\le2\|Y\|_\infty\|Z\|_\infty$
for the terms with $I_i$ leaves, after division by $x^{-2}M_1M_2$,
at most
\[
 C\big[(\log x)^{21p-q}+(\log x)^{18p}x^{-b}\big].
\]
Indeed the two support counts cost $(\log x)^{16p}$ and the
background bound costs another $(\log x)^{2p}$. The square of the
irrelevant bracket is absorbed by the bracket for large $x$.
The right side tends to zero. Substitute these estimates in
\eqref{fp:eq-exact-covariance}.
\end{proof}

\subsection{Comparison of the clean retained law}
\label{fp:sec-clean-kernels}

Let $u=a/a_s$. In the admitted small logarithm chart write the
background plaquette on the retained datum at level $j$ as
$\exp(iu^2Y_{j,p})$. Thus $Y_{j,p}$ is its traceless Hermitian
curvature coordinate in units of $u^2$, with the half-trace norm
of the physical setting. For $j=j_0$ it is the
background logarithm in $X_i$. For each plaquette pair use one
common colour coordinate for its readouts throughout the
comparison. The held aligned pair at level $j$, modulo a
simultaneous colour rotation, is the same measurable function of
the retained state in both adjacent steps. One may choose any
measurable representative; each box comparison may rotate both
reads together because the polynomials used below are invariant.
A simultaneous colour rotation preserves the scalar colour covariance in
\eqref{fp:eq-slice-covariance}. Independent rotations of the two
reads require the corresponding transport in that covariance.
Individual curvature squares are gauge invariant.

Here are the inputs from the reference-window comparison. They
concern the complete states of the original run on the reference
torus. Write \(\lambda_j\) for the unsourced standing law with
the block field and stored slots, and \(\Pi_j\) for the positive
joint law of the original step from \(\lambda_{j-1}\) to
\(\lambda_j\). Its marginals are these two laws.
Use bounded measurable extensions of the curvature reads to all
states, with the same read in both adjacent comparisons. At the
witness set the curvature coordinate to zero whenever its
passed-test indicator fails. Its square then carries exactly
the indicator of \(X_i\). Require the extensions to obey
\[
             |Y_{j,p}|\le C g_x P
\]
on every retained state, including every excluded part, with the
common factor \(P\) fixed below. On the witness passed-test set
this bound must hold for the actual read, so the extension does
not change \(X_i\). Outside admitted reads an extension can be
clipped or set to zero. Boundedness without the factor \(g_x\)
would not control the excluded mass relative to the signal.

At step \(j\), the admitted part of \(\Pi_j\) is selected by a
good event measurable in the held exterior, after summing the
standing parts. The selector is constant on the product window;
its complement has mass \(p_j\). On the good exterior the
normalized box law below is the disintegration of this original
law. The same box integral normalizes the conditional measure
and occurs in the exterior mixture weights. The read set
contains every pivot and contour dependency of the pair. A
restriction on an already deleted slot needs its own comparison
and is not imposed by this selector. The argument uses the
original adjacent laws without conditioning the entire path on
all good events simultaneously.

For each pair of plaquettes and each step let $S$ be this finite
set of free coordinates. The following bounds hold
uniformly over the admitted exterior and all such pairs. A common
factor $P=(\log x)^B\ge1$ bounds the witness room, read dimension,
maximum absolute free coordinate on the window, scaled curvature $g_x^{-1}|Y_{j,p}|$, and the
polynomial chart constants used below. Increasing the fixed
exponent $B$ is allowed.

\begin{enumerate}
\item At fixed exterior the normalized actual marginal on $S$
has density ratio in $[e^{-\xi_j},e^{\xi_j}]$ with respect to the
same normalized flat box-window marginal. Here $\xi_j\le1$.
Both flat Gaussians before restricting any window are centred
in the same free coordinates. The local box Gaussian has
window-tail probability at most $t_j\le1/2$. Its unwindowed
covariance compressed to $S$ differs in Hilbert--Schmidt norm by
at most $d_j$ from the corresponding full flat torus marginal.
Both compressed covariances have eigenvalues in one fixed
positive interval. This comparison includes the coarea and
gauge densities, the full interaction, the exterior force and
the cut. The marginal before removing the cut is on the same
coordinates as the one used in the readout.

\item On each selected full product window, the bounded
previous-read extension agrees with a curvature readout
\(T_{j,p}(y;e)\), and the passed-test indicators in that
extension equal one. Here \(y\) consists of the free coordinates
on \(S\) and \(e\) is the held exterior. The retained read
\(Y_{j,p}\) is held there. The direct readout estimate is
\begin{equation}\label{fp:eq-direct-readout}
 \sup_{y\,\mathrm{in\ window}}
 |T_{j,p}(y;e)-Y_{j,p}-B^{\rm loc}_{j,p}y|
       \le Cg_xP^4(\epsilon_j+g_x).
\end{equation}
The linear row $B^{\rm loc}_{j,p}$ is the
flat row for this local read, using the actual physical
weight of the step. Embed it by zero into the full flat
coordinate set. There is a full flat row $B_{j,p}$ on that
set with
\[
 \|B^{\rm loc}_{j,p}-B_{j,p}\|_{2\to2}\le Cg_xP\epsilon_j,
 \qquad \|B_{j,p}\|_{2\to2}
             \le Cg_x L^{-2(j-1-j_0)} .
\]
Choose the full centred flat comparison law
\(\gamma_j^0\) to be the zero-next-datum soft-contour Gaussian
of \cite[Proposition~3.3]{GI}, with its fixed positive contour
weight. Its free-coordinate covariance is uniformly bounded
on the admitted relative weight range. Write
\(\widehat g_{j-1}=\theta_{j-1}^{-1/2}\) and
\(\tau_{j-1}=x_{j-1}/\theta_{j-1}\).
The full row \(B_{j,p}\) is \(u^{-2}\widehat g_{j-1}\) times
its linear lifted-curvature row. The local comparison must use
this same flat coordinate system and row. By
\cite[Propositions~3.3 and~4.3]{GI}, scaling the action by
\(\tau_{j-1}\) while keeping the contour weight fixed scales
its curvature covariance by \(\tau_{j-1}^{-1}\). Consequently
\begin{equation}\label{fp:eq-slice-covariance}
 \mathbb E_{\gamma_j^0}
       [(B_{j,p}y)\otimes(B_{j,q}y)]
 =u^{-4} x_{j-1}^{-1}(P_{j-1}-P_j)(p,q)\,I_3.
\end{equation}
Here \(\mathcal C_{(j)}=P_{j-1}-P_j\) is the unit raw-flux
slice. The physical factor is exactly
\(\widehat g_{j-1}^2\tau_{j-1}^{-1}=x_{j-1}^{-1}\).

\item On the admitted terminal part,
\(|Y_{j_1,p}|\le Cg_xP L^{-2n}\). Its complement has mass
\(p_{\rm top}\), and the bounded extension is used there.
The choices of radii, localization buffers and excluded masses give
\begin{equation}\label{fp:eq-clean-smallness}
 P^{50}\Delta\longrightarrow0,\qquad
 \Delta:=L^{-2n}+n g_x+p_{\rm top}+
       \sum_{j_0<j\le j_1}
        (p_j+\xi_j+\sqrt{t_j}+d_j+\epsilon_j).
\end{equation}
The fixed power $50$ is a sufficient reserve for the finite
degree comparison below. The geometry and real window widths
fix $B$ before the analytic decay and window-tail reserves are
required. For example, the excluded masses are
superpolynomially small if their events are covered by at most
a fixed power of \(x\) local contacts throughout the segment,
each of probability at most \(Ce^{-c(\log x)^p}\), with \(p>1\).
This union bound concerns the original adjacent laws and their
actual determining sets.
\end{enumerate}

One sufficient source of \eqref{fp:eq-direct-readout} is a
product window centred at zero, with \(\|y\|_2\le P^{3/2}\),
and the estimates
\[
 \begin{split}
 |T_{j,p}(0;e)-Y_{j,p}|&\le Cg_xP\epsilon_j,\\
 \|D_yT_{j,p}(0;e)-B^{\rm loc}_{j,p}\|&\le Cg_xP\epsilon_j,\\
 \sup_{y\,\mathrm{in\ window}}\|D_y^2T_{j,p}(y;e)\|
     &\le Cg_x^2P.
 \end{split}
\]
Taylor's formula along \([0,y]\) then bounds the constant,
linear and quadratic errors by
\(Cg_xP\epsilon_j\), \(Cg_xP^{5/2}\epsilon_j\) and
\(Cg_x^2P^4\), respectively. The direct readout estimate
can instead be supplied without these derivative hypotheses.

The normalized density-ratio estimate also has a precise
sufficient condition. At fixed exterior, suppose the actual
conditional density on the same cut is proportional to
\(e^{-V}\) times the flat reference probability, where \(V\)
is real and its oscillation is at most \(\xi_j\).
Its normalizing integral lies between
\(e^{-\sup V}\) and \(e^{-\inf V}\), so the normalized
ratio lies in \([e^{-\xi_j},e^{\xi_j}]\).
Conditional expectation gives the same bounds for the marginal
on \(S\). Establishing this oscillation requires the complete
conditional interaction and exterior force on that cut.
A bound on only its quadratic part does not suffice.

\begin{lemma}[Transport of polynomial readouts]
\label{fp:lem-retained-covariance}
Under the adjacent-law inputs above, take two read plaquettes
\(p,q\) and a homogeneous polynomial \(H\) of degree
\(m\in\{0,2,4\}\) in their six scaled curvature coordinates.
Its sum of absolute coefficients is at most one, and it is
invariant under simultaneous \(SO(3)\) rotation of the two reads.
Put $W_j=(Y_{j,p},Y_{j,q})$ and let $Z_j$ have the law of the
six-dimensional vector of full flat readouts of $\gamma_j^0$,
with the $Z_j$ independent over levels. For all sufficiently
large \(x\), uniformly on the stated input family,
\begin{equation}\label{fp:eq-polynomial-comparison}
 \left|\mathbb E_{\lambda_{j_0}} H(W_{j_0})
        -\mathbb E H\left(\sum_{j_0<j\le j_1}Z_j\right)\right|
 \le Cg_x^m P^{20}\Delta .
\end{equation}
A constant polynomial has zero error.
\end{lemma}

\begin{proof}
First, at a given fixed exterior replace the actual read by
$W_j+(B^{\rm loc}_{j,p}y,B^{\rm loc}_{j,q}y)$. The direct
estimate \eqref{fp:eq-direct-readout} bounds this replacement
error by \(Cg_x P^4(\epsilon_j+g_x)\). Increase the floor for
\(x\) until every \(\epsilon_j\le1\), as permitted by
\eqref{fp:eq-clean-smallness}. The retained read and the
linear Gaussian increments on that window have size at most
$Cg_xP^3$, increasing fixed constants if necessary.

Define polynomials
\[
 H_j(w)=\mathbb E H\left(w+\sum_{j_0<l\le j} Z_l\right),
 \qquad H_{j_0}=H .
\]
Every slice has scalar colour covariance in the common pair
frame. Its law, and therefore each \(H_j\), is invariant under
a simultaneous rotation of the two reads. Expansion of \(H\)
into monomials bounds each derivative and moment norm by the
sum of its individual bounds.
The bound on the full rows and the geometric sum of their
squared sizes give bounded moments of each coordinate of the
Gaussian sum, of order $O(g_x^b)$ for a moment of degree $b$.
This expansion bounds its derivative of order $r$
at $|w|\le Cg_xP^4$ by $C g_x^{m-r}P^{4(m-r)}$.
This gives a local Taylor replacement error bounded by
$C g_x^m P^{20}(\epsilon_j+g_x)$ for $H_{j-1}$.

The density ratio in the first input changes a bounded
window expectation by at most $C\xi_j$ times its supremum.
To remove the window, write the full expectation as the two
integrals on the window and its complement. Cauchy--Schwarz
bounds the latter by $\sqrt{t_j}$ times the $L^2$ norm of the
polynomial. The same Gaussian moment estimates, with the held
$W_j$ bounded by $Cg_xP$, bound this norm by $Cg_x^mP^{20}$.
The normalization of the window costs at most a factor two.
For the difference of unwindowed compressed Gaussians interpolate
their covariance matrices linearly. If $C_v$ is the interpolant
and $D$ the covariance difference, the derivative of its density
is that density times
\[
 \tfrac12\big(y^{\mathsf T}C_v^{-1}DC_v^{-1}y
                       -\operatorname{tr}(C_v^{-1}D)\big).
\]
The second moment of the last factor is
$\tfrac12\|C_v^{-1/2}DC_v^{-1/2}\|_{\rm HS}^2$.
Uniform bounds on the compressed covariances and
Cauchy--Schwarz against the polynomial therefore give error
$Cg_x^mP^{20}d_j$. On the full flat torus now use the same
Gaussian to replace the local row by the full row. Its linear
read difference has every fixed $L^b$ norm at most
$C_b g_xP\epsilon_j$. The polynomial identity for differences
and Gaussian moments give the same cost with $\epsilon_j$.

On the good exterior the estimates compare
\(H_{j-1}(W_{j-1})\) under the original box law with
\(H_j(W_j)\) at the held exterior. Their uniformity permits
integration against the original exterior weights. On the
excluded part, bounded extensions and Gaussian moments bound
both quantities by \(Cg_x^mP^{20}\), so its cost is at most
\(Cg_x^mP^{20}p_j\). The marginal identities give
\[
 \begin{split}
 &\mathbb E_{\lambda_{j_0}}H(W_{j_0})
       -\mathbb E_{\lambda_{j_1}}H_{j_1}(W_{j_1})\\
 &\qquad=\sum_{j_0<j\le j_1}
   \mathbb E_{\Pi_j}\bigl[H_{j-1}(W_{j-1})-H_j(W_j)\bigr].
 \end{split}
\]
A carried observable is the Radon--Nikodym derivative of its
weighted pushforward with respect to the uninserted pushforward,
as one checks against bounded functions of the output. At a
merge both measures are sums with the original weights. Thus
this identity preserves the exterior distribution. At the top,
compare \(H_{j_1}(W_{j_1})\) with \(H_{j_1}(0)\) on the admitted
part and use the bounded extension on its complement. The cost
is \(Cg_x^mP^{20}(L^{-2n}+p_{\rm top})\). Summing proves
\eqref{fp:eq-polynomial-comparison}. Comparing full moments
before forming their covariance retains the covariance of
conditional means.
\end{proof}

\begin{proposition}[The reference covariance]
\label{fp:prop-clean-reference}
Suppose the witness and clean-window inputs of
\cref{fp:lem-witness-errors,fp:lem-retained-covariance} hold on
one reference torus. Use the admitted flat hierarchy of
\cref{fp:lem-free-kernel-bounds} through the two stops.
Assume the enlarged supports are in its separated kernel
range at the witness, and $x_{j-1}=x+O(n+1)$ on the segment.
Then
\begin{equation}\label{fp:eq-reference-wilson}
 S_{2;\star}^T(f_1,f_2)
  =\frac32\zeta^2 x^{-2}
       \sum_{p,q}\widetilde f_1(p)\widetilde f_2(q)
                           G_{j_0}(p,q)^2
                +o(x^{-2}) M_1M_2 .
\end{equation}
The error is uniform in the admitted bare data and positions.
\end{proposition}

\begin{proof}
Let $G_{j_0}=C^{\rm bg}_{K,s}$ and $G_{j_1}$ be its nested
coarser kernel. By \cref{fp:lem-free-kernel-bounds},
\[
 |G_{j_0}(p,q)|\le C a^4 d^{-4},\qquad
 |G_{j_1}(p,q)|\le C u^4 L^{-4n}.
\]
On the same enlarged supports each unit slice is bounded by
\[
 |\mathcal C_{(j)}(p,q)|\le C u^4 L^{-4(j-1-j_0)}.
\]
With the bounded witness extension, the passed-test Wilson
read is \(\tfrac12u^4|Y_{j_0,p}|^2+O(u^8|Y_{j_0,p}|^4)\).
This expansion holds on the witness law, with both sides zero
when the test fails. No global clean selection is removed.
There are at most $C\widehat d^4u^{-4}$ plaquettes in
each read support, including its fixed enlargement.
Since $|Y_{j_0,p}|\le Cg_xP$ and $\widehat d\le P$, the
Wilson Taylor error is at most $Cg_x^4P^8M_i$ in the one-point
read. Each quadratic read is bounded by $Cg_x^2P^6M_i$ on the
entire witness law, using the bounded extension. The resulting covariance error is at most
$C(g_x^6P^{14}+g_x^8P^{16})M_1M_2=o(x^{-2})M_1M_2$.
Apply \cref{fp:lem-retained-covariance} to the square of one
curvature read and the product of squares at two plaquettes.
Use \(\Delta\) from \eqref{fp:eq-clean-smallness}.
After smearing, the factors $u^4$ and $u^8$ cancel the
plaquette counts. The degree-two mean errors are at most
$C g_x^2P^{24}\Delta M_i$, and the degree-four product error
is at most $C g_x^4P^{28}\Delta M_1M_2$. Each Gaussian mean is
bounded by $C g_x^2P^4 M_i$, while the actual quadratic mean
has the $P^6$ bound just established. Subtracting the products
of means thus costs at most $C g_x^4P^{30}\Delta M_1M_2$.
The reserve $P^{50}\Delta\to0$ pays each of these losses.
The polynomials \(|Y_p|^2\) and \(|Y_p|^2|Y_q|^2\)
are invariant under simultaneous colour rotations. Their
coefficient norms are \(3\) and \(9\); rescaling them to the
unit coefficient ball only changes the fixed constants.
Thus the argument uses degree-two and degree-four invariant
polynomials on a pair of reads.

For the sum of independent slices, Wick's formula gives
the covariance of two half-squares as $3/2$ times the
square of the one-colour covariance. Thus its kernel is
$\tfrac32[\sum_j x_{j-1}^{-1}\mathcal C_{(j)}(p,q)]^2$.
In particular the square contains the cross terms between
different steps. They would be lost by summing only the
one-step covariances at a zero retained field. Replacing
the weights $x_{j-1}^{-1}$ by $x^{-1}$ has kernel error
$O(x^{-2}(n+1)u^4)$ before squaring, by the geometric slice
bound. The full square error after smearing is
$O(x^{-3}(n+1)\widehat d^8)M_1M_2=o(x^{-2})M_1M_2$.
The slice sum now telescopes to $G_{j_0}-G_{j_1}$.
The cross product with $G_{j_1}$ and its square have
total smeared cost at most
\[
 Cx^{-2}\big[(\widehat d L^{-n})^4
              +(\widehat d L^{-n})^8\big]M_1M_2 .
\]
Here $n=\lceil2\log_L\bar x_s\rceil$ and $\widehat d$ is
polylogarithmic in $x$, so this is $o(x^{-2})M_1M_2$.
Equation \eqref{fp:eq-witness-errors} proves the result.
If a local computation uses witness-lattice translates, the
bare covariance and the free kernel have the same value on
these translates. One may make all the replacements after
averaging them; no finer translation invariance of the block
fields is used.
\end{proof}

\section{A common normalization with a finite reference torus}
\label{fp:sec-finite-reference-transfer}

We keep the bare datum fixed while changing volume. Write
$\ell_\star$ for the reference side and assume every compared
torus side is an integer multiple of it. Choose an allowed
board side $l_c=L^m$ with $\ell_\star/l_c$ even and with
room for a unit ball at fixed positive distance from its
faces. This side can be fixed at the reference coupling.
The periodic array of sources below is obtained by the
same cell reflections on all of these tori. For a physical test
$f$ use $\|f\|_\sharp=\max(\|f\|_\infty,c_M\operatorname{Lip}f)$,
with the fixed positive constant $c_M$ of the SI source norm.

\begin{lemma}[Replication of a reflected source]
\label{fp:lem-replicated-partition}
Let $w$ be a complex source in the interior of one board cell,
with no sourced plaquette crossing a cell face. Form the
reflected array $W_\ell$, conjugating the source under an odd
number of cell reflections. Then the Wilson partition integrals
with Haar probability measures satisfy
\begin{equation}\label{fp:eq-partition-replication}
 0\le Z_\ell(W_\ell)
       \le Z_{\ell_\star}(W_{\ell_\star})^{(\ell/\ell_\star)^4}.
\end{equation}
The unsourced bare inverse coupling is the same on every torus.
\end{lemma}

\begin{proof}
First lengthen one side, keeping the other three fixed. Use the
group of spatial bond variables on a slice as the integration
space, with product Haar probability. Between two slices at a
cell face there is no source. If their spatial bonds are $u,v$,
the temporal bonds between slices give, on integration, the
kernel obtained from
$\prod_b c(u_b (v_b^h)^{-1})$ by averaging the slice gauge
$h$, where $c(g)=e^{-\beta(1-\operatorname{Tr}g/2)}$ and
$v_b^h$ is the usual site gauge action on the bond $b$.
This is the crossing-plaquette product after cyclically
permuting its traces. Thus the crossing operator is $K=CP$,
where $C$ is convolution by the product of $c$ over slice bonds
and $P$ is Haar averaging over slice gauges. Centrality of $c$
makes the convolution commute with the left and right group
actions on each bond, hence with $P$.

The expansion
$c=e^{-\beta}\sum_{n\ge0}(\beta/2)^n\chi_{1/2}^n/n!$
has nonnegative character coefficients. Indeed the powers
are characters of tensor powers of the defining representation.
At the identity the sum of coefficients times the representation
dimensions is one. For one bond a coefficient $a_\lambda$ in
dimension $d_\lambda$ gives convolution eigenvalue
$a_\lambda/d_\lambda$ with multiplicity $d_\lambda^2$.
Consequently $C$ is positive trace class on the finite product
slice space. The commuting orthogonal projection $P$ makes
$K$ positive trace class as well. Its square root $S$ is
Hilbert--Schmidt.

Let $B$ be the transfer through one slab between two board
faces with the crossing factors at its ends removed. It is the
product of the spatial multiplication factors and the integrated
interior temporal kernels. These are bounded operators since the
slice space is a compact product with finite Haar measure and
the profile is bounded. The next slab is the reflected
conjugate slab and has operator $B^*$. This follows by reversing
the order of the slices and of the temporal bond orientations,
which preserves Haar measure and the plaquette traces. Any
transverse board array is held fixed in this argument. With
$A=SBS$, the partition integral over a whole period is
$\operatorname{Tr}(AA^*)^h$ if the period has $2h$ slabs.
Indeed before splitting each crossing factor it is the
cyclic integral of $KB,KB^*,\ldots$. Each of these factors has a
Hilbert--Schmidt kernel. The trace identity for a product of two
Hilbert--Schmidt kernels, with the other factors composed by
Fubini, gives that cyclic integral. Cyclicity of the trace then
splits the crossing factors into the two square roots. This
also covers a slab consisting of a single time slice, when $B$
is a multiplication operator.

The positive trace-class operator $T=(AA^*)^h$ has
$\operatorname{Tr}T^b\le(\operatorname{Tr}T)^b$ for every
positive integer $b$, by its nonnegative eigenvalues. Repeating
the period $b$ times gives exactly the array on the longer side.
This proves the one-direction inequality. The number of cells
is even in every direction, so successive replication in all
four directions preserves the same array. The intermediate
rectangular tori use only finite Wilson integrals. They require
no effective-action construction. This proves
\eqref{fp:eq-partition-replication}.
\end{proof}

\begin{lemma}[Comparison of vacuum densities]
\label{fp:lem-vacuum-density}
Let $b_j=L^{j-K}$, $0\le j\le K$. Suppose the once-counted
unsourced vacuum on each compared torus has the form
\[
 E_{\ell,K}=\sum_{j=0}^K\left[
 (\ell/b_j)^4 c_{j,K}
 +\sum_{z\in (b_j\mathbb Z/\ell\mathbb Z)^4}
                \sum_{D\ni z}v_{\ell,j}(D,z)\right].
\]
The numerical coefficients $c_{j,K}$ and the fixed auxiliary
data are common. There are rooted infinite companions
$v_{\infty,j}$ with the same translation-covariant sums per
root. Suppose their rooted domain counts at length $m$ are
bounded by $C e^{\kappa_D m}$, and their individual absolute
values by $C e^{-\kappa_E m}$ with $\kappa_E>\kappa_D$.
The same bounds hold for the finite domains. At each root
assume a bijection of finite and infinite rooted reads of length
$m<\ell/(C_0b_j)$ which preserves their vacuum values.
Then for $\ell\ge\ell_0>0$ there is a common number $A_K$
such that
\begin{equation}\label{fp:eq-vacuum-density}
 \left|\ell^{-4}E_{\ell,K}-A_K\right|
 \le C\sum_{j=0}^K b_j^{-4}e^{-c\ell/b_j}
 \le C_{\ell_0}e^{-c\ell/2}.
\end{equation}
The constants are independent of the cutoff. The number $A_K$
is allowed to depend on the cutoff and the common bare datum.
\end{lemma}

\begin{proof}
The companion sum $\bar v_j$ at one root is absolutely
convergent. Set $A_K=\sum_j b_j^{-4}(c_{j,K}+\bar v_j)$.
Short rooted values cancel in the difference. The two tails
are bounded by geometric sums with the positive gap
$\kappa_E-\kappa_D$. Dividing by the physical volume yields
the first inequality. For the second, put $u=K-j$ and bound
\[
 \sum_{u\ge0} L^{4u}e^{-c\ell L^u}
 \le e^{-c\ell/2}
              \sum_{u\ge0}L^{4u}e^{-c\ell_0L^u/2}.
\]
The last sum converges. Thus $A_K$ cancels exactly in a
comparison with the reference torus and need not converge.
In the scalar convention of \cite[Proposition~2.5]{RT} the numerical free-bond and
gauge factors go in $c_{j,K}$. One can point an unpointed
remainder by allocating its value equally to its roots.
The assumed equality is an equality of the actual rooted
vacuum values on nonwrapping complete reads. Equality of
bare weights alone does not imply it.
\end{proof}

\fpneedspace{24\baselineskip}
\begin{proposition}[One analytic cell subtraction for both measures]
\label{fp:prop-shared-normalizer}
Use the comparison geometry and the hypotheses of
\cref{fp:lem-vacuum-density}. Suppose the actual reference
torus has the ordinary complex source upper bound of \cite[Proposition~8.1]{SI}
\[
 |Z_{\ell_\star}(W)|
  \le \exp\{E_{\ell_\star,K}+\operatorname{Re}\mathcal E_\star(W)
                    +C_+\ell_\star^4\},
\]
and every compared torus, including the reference, has the
unsourced floor
$Z_\ell(0)\ge \exp\{E_{\ell,K}-C_-\ell^4\}$.
Assume the constants are uniform at the fixed reference
coupling. In the reference source number, each domain term
is holomorphic in the source on the ordinary radius $2r$,
vanishes at zero and depends only on its complete source
read. At scale $b_j$ root each term at a level site in its
complete read. At length $m$ assume the absolute bound
$C e^{-\kappa_E m}$ on the ordinary source ball. Require the
complete read together with its root to have diameter at most
$C b_j(m+1)$ and the rooted count to be at most $C e^{\kappa_D m}$,
with $\kappa_E>\kappa_D$.
Separate point terms, if used, are holomorphic on the same
ordinary source ball, vanish at zero and have only their own
point read.

For any real Lipschitz tests $F_1,F_2$ supported in one unit
ball, there is one holomorphic function $E^\circ(z)$,
$E^\circ(0)=0$, for which
\begin{equation}\label{fp:eq-shared-normalizer}
 \left|e^{-E^\circ(z)}
   \omega(e^{z_1\mathcal O_a(F_1)+z_2\mathcal O_a(F_2)})\right|
     \le Q
\end{equation}
on $\sum_i |z_i|\|F_i\|_\sharp<r$.
Here $\omega$ can be the reference measure, any compared
torus measure or any periodic Wilson Gibbs state at this
datum. For the last assertion assume the compared tori
contain a compatible cofinal sequence. The constants are
uniform in the cutoff, volume, position and states.
\end{proposition}

\begin{proof}
Place the common support ball inside the chosen cell by a
joint bare lattice shift. Translation invariance lets us
return to the original position. Use the cell source
$w_z=z_1F_1+z_2F_2$ and its reflected conjugate array $W_\ell$.
The even and odd profiles separately vanish near faces, so
joining copies does not increase their Lipschitz constants.
Using the two sets of coefficients $z_i$ and $\bar z_i$ gives
ordinary source size at most $2\sum_i|z_i|\|F_i\|_\sharp$.
The reflected dissemination of $e^{\mathcal O_a(w_z)}$
is $e^{\mathcal O_a(W_\ell)}$, since each sourced plaquette
is internal and reflection preserves its trace. With
$N_\ell=(\ell/l_c)^4$, its finite chessboard seminorm is
$[Z_\ell(W_\ell)/Z_\ell(0)]^{1/N_\ell}$. The numerator is
nonnegative. By \cref{fp:lem-replicated-partition} and the
two partition bounds, this seminorm is at most
\[
 \exp\left\{
 N_{\ell_\star}^{-1}\operatorname{Re}
     \mathcal E_\star(W_{\ell_\star})+C l_c^4
 \right\}.
\]
Here we used \eqref{fp:eq-vacuum-density} for the difference
of the two vacuum densities. The common divergent number
has cancelled.

Define $E^\circ$ by keeping the reference source terms whose
complete read lies inside one open board cell and averaging
over the cells. Keep every nonzero point term. Use the value
of each even term and the conjugate of each odd term. The
odd actual value is $h(\bar z)$ for a holomorphic $h$, so
$\overline{h(\bar z)}$ is holomorphic and has the same real
part. All the kept terms are evaluated on the source in
their own cell. They vanish at zero. Normal convergence
on smaller source discs follows from the rooted bounds
at the fixed cutoff. This defines the same $E^\circ$ for all
the measures in the statement.

A discarded nonzero term meets a source support and reads
outside that support's cell. Its diameter is at least
the fixed clearance. At scale $b_j$ its length is at least
$c/b_j-C$. Charge the term to one cell in which it reads a
nonzero source. Its root is within $C b_j(m+1)$ of the support
ball in that cell, so there are at most
$C(1+b_j^{-1})^4(m+1)^4$ possible roots. The strict entropy
margin absorbs the polynomial in $m$. Summing over the
allowed lengths and then the scales, the omitted absolute
source number per cell is bounded by
\[
 C\sum_{u\ge0}(1+L^u)^4 e^{-c'L^u},
 \qquad u=K-j.
\]
The constant incorporates the fixed localization side and collar. Thus the finite
seminorm after subtracting $E^\circ$ is bounded by a
common $Q$. The one-cell chessboard inequality gives
\eqref{fp:eq-shared-normalizer} on the tori. The fixed-spacing
all-state board bound of \cite[Theorem~C]{LP} gives it along
the same cofinal family for every periodic state. This
argument uses bounded cylinder exponentials at the fixed
cutoff. In particular it does not select the state as a
weak limit of these particular tori.
\end{proof}

\begin{lemma}[Transfer of a local covariance]
\label{fp:lem-local-transfer}
Let two translation-invariant laws at the same bare
spacing, with the same linear source convention, obey
\eqref{fp:eq-shared-normalizer} for every unit-ball test pair,
with a common $E^\circ$ for the two laws. The tests in this
assumption may overlap and have either sign. For a pair in
\cref{fp:def-shape}, with $a\le d$ sufficiently small, the
difference of their connected covariances satisfies
\begin{equation}\label{fp:eq-local-transfer}
 |S_{2;1}^T(f_1,f_2)-S_{2;2}^T(f_1,f_2)|
        \le C_\vartheta d^2 M_1M_2 .
\end{equation}
For tori, take supports and the shifted unit ball in one
common lift.
\end{lemma}

\begin{proof}
For any unit-ball pair, apply the local residual Cauchy bound
of \cite{SI} to $e^{-E^\circ}$ times the two generating
functions. Their values at zero are one. Equivalently the
bound $Q$ and radial Schwarz give a logarithm bounded by
$\log 2$ on a bidisc whose coordinate radii are
$r/[4(Q+1)\|F_i\|_\sharp]$. The two mixed derivatives of
$E^\circ$ cancel. If $D$ is the difference of the two
covariances, then
$|D(F_1,F_2)|\le C\|F_1\|_\sharp\|F_2\|_\sharp$.
A zero test gives zero directly.

The two supports lie in one radius-$5d$ ball. Pack
$J\ge c d^{-4}$ disjoint lattice translates of that ball
inside a unit ball, using shifts on a grid with a lattice
spacing comparable to $d$. For independent signs
$\varepsilon_v$ indexed by the shifts put
$F_i=\sum_v\varepsilon_v\tau_v f_i$. Both the supremum
and Lipschitz norms are unchanged or decreased, because the
supports of the copies are disjoint and each test vanishes
at its ball boundary. For points in different balls split
their connecting segment at the two boundaries to obtain
the Lipschitz assertion. The sign average of $D(F_1,F_2)$
is $J D(f_1,f_2)$ by bilinearity and joint translation
invariance. Divide by $J$ and use
$\|f_i\|_\sharp\le C_\vartheta d^{-1}M_i$.
This gives \eqref{fp:eq-local-transfer}. Separate analytic
normalizers would leave their mixed derivative difference
in this calculation.
\end{proof}

\begin{theorem}[The short-distance two-point law]
\label{fp:thm-two-point}
Use the witness and running source of
\cref{fp:prop-witness,fp:prop-source-normalization} at one
fixed sufficiently small reference coupling, the actual
reference comparison of \cref{fp:prop-clean-reference},
and the common finite-reference inputs of
\cref{fp:prop-shared-normalizer}. Fix the common integer
\(n_0\ge1\) above which those formation and reference inputs
are supplied, and require \(K\ge s(d)+n_0\).
Let \((f,h)\) belong to Definition~\ref{fp:def-shape} at
separation \(d\). Then below a common
$d_0>0$, on every admitted torus and in every admitted
periodic Gibbs state,
\begin{equation}\label{fp:eq-two-point-law}
 S_{2;a,\rho}^T(f,h)=\frac9{\pi^4}\,\zeta_K(d)^2g_K(d)^4
       \mathcal K(f,h)(1+r_{K,\rho}(f,h)),
 \qquad |r_{K,\rho}(f,h)|\le\tfrac12 .
\end{equation}
Here \(\rho\) denotes the compared finite-volume law or periodic
state; the remainder may depend on it. The error tends to zero
as \(d\downarrow0\), uniformly over
the bare data, support positions and compared measures
for which the displayed input estimates are uniform.
The two factors $\zeta_K(d)$ and $g_K(d)$ are those of the
same run at the same bare datum throughout the comparison.
\end{theorem}

\begin{proof}
To put the same pair on the reference torus, lift it to the
common radius-\(5d\) ball and extend each test by zero outside
its supporting ball. Stopping a line segment at its first
boundary point, where the test vanishes, preserves the
Lipschitz bound. Periodization into any torus of side at least
\(40d\) preserves that bound, the masses, support gap and
\(\mathcal K\); split shortest segments at boundaries of
the disjoint lifted copies.
For the actual signed profiles $\widetilde f_i$, the tent estimates give
$\|\widetilde f_i-f_i\|_\infty=o(1)M_i$ and an enlargement
by at most $CMa_s$. The supports remain separated by at
least $d/2$ below a common threshold, by the witness room.
For small enough \(d\), the witness bounds \(a\le a_s\) and
\(d/a_s\to\infty\) give \(10a\le d/2\), \(a_s\le d/2\) and
\(d\le\ell_\star/200\). The enlarged supports then stay in
the separated kernel range. There are at least two sites in
each direction at the used coarse levels.
Apply \cref{fp:prop-clean-reference} with these signed profiles.
The separated free background bound
$|G_{j_0}(p,q)|\le C a^4d^{-4}$ and the two plaquette
counts show that replacing the signed profiles by the
originals costs $o(x^{-2})M_1M_2$. This uses an absolute
estimate before any division by a positive main term.
The hierarchy and distance hypotheses of
\cref{fp:prop-free-reference} hold for the original pair.
That proposition replaces the unit raw-flux square sum by
$6\pi^{-4}\mathcal K(f,h)(1+o(1))$.
The factor $3/2$ in \eqref{fp:eq-reference-wilson} then
gives $9\pi^{-4}$.

The shape estimate and running-source lower bound give
\[
 \frac9{\pi^4}\zeta_K(d)^2 x^{-2}\mathcal K(f,h)
      \ge c_\vartheta x^{-2}\|f\|_\infty\|h\|_\infty.
\]
The relative reference error therefore tends to zero.
The free replacement also gives
\[
 \sum_{p,q}f(x_p)h(x_q)G_{j_0}(p,q)^2\ge
                         c_\vartheta M_fM_h .
\]
Thus the reference covariance equals this unmodified lattice sum times
\(3\zeta^2x^{-2}/2\) with relative error at most \(1/4\),
after decreasing the common threshold.
Apply \cref{fp:lem-local-transfer} with the finite reference
and the given measure. Its relative error is
$C_\vartheta d^2x^2$, which tends uniformly to zero since
$x\asymp\log(1/d)$. The common-normalizer argument uses
the same physical source with the same bare datum, without
changing the Wilson coefficient or rescaling the insertion.
Decrease $d_0$ until the total relative error is at most
$1/2$. This proves \eqref{fp:eq-two-point-law}.
\end{proof}


\fpneedspace{18\baselineskip}
\begin{corollary}[Transfer with a cofinal reference]
\label{fp:cor-cofinal-reference}
For each fixed bare datum suppose there is one compatible
cofinal sequence of finite tori on which
\eqref{fp:eq-reference-wilson} holds for \(K\ge s(d)+n_0\),
with one fixed \(n_0\ge1\) and a common positive distance
threshold for this cofinal comparison. Bound its additive remainder by
\(\epsilon_{\rm ref}(d)x^{-2}M_1M_2\), where one
\(\epsilon_{\rm ref}(d)\to0\) works uniformly over the cofinal
sides above a common positive minimum, all admitted bare data,
shapes and support positions. Use the same actual profiles and
hierarchy as in \cref{fp:prop-clean-reference}. The sequence at one datum
works for every fixed admissible pair there with the same
witness and coefficient factors. Suppose the fixed-spacing
state-transfer hypotheses of \cite[Theorem A]{SI} hold uniformly, with
a common cell subtraction for all periodic states at the datum.
Then the conclusion of \cref{fp:thm-two-point} holds for the
periodic states. It needs no comparison of vacuum densities
with one fixed finite torus in this case.
\end{corollary}

\begin{proof}
First replace the pointed profiles by the originals using their
absolute tent error and the separated kernel bound. Apply
\cref{fp:prop-free-reference} on each torus; its error is uniform
once the side is above the common minimum. Dividing the resulting
additive error by the positive shape floor gives the same
relative reference law as in the finite-reference proof.
At fixed spacing take a local weak limit of the cofinal torus
laws, using compactness and the finite-range Gibbs specification
as in \cite{LP}. This gives a periodic Wilson Gibbs state at
the same bare datum. The covariance of the fixed compact
tests is built from bounded continuous cylinder functions
at that spacing. The uniform reference estimate therefore
passes to this state, with the same $\zeta_K(d)$ and $g_K(d)$.
The common cell subtraction of \cite[Theorem A]{SI} serves the limit state and
any periodic state at that datum. Apply
\cref{fp:lem-local-transfer} and the positive shape lower
bound exactly as in the finite-reference proof. Uniformity
over the cofinal sequence for fixed pairs is used before
taking the volume limit. A distance threshold tending to
zero along that sequence would not suffice for this passage.
\end{proof}


\fpneedspace{23\baselineskip}
\section{Ultraviolet bounds on the prescribed class}
\label{fp:sec-consequences}

\begin{corollary}[Bounds and conditional passage to a limit]
\label{fp:cor-ultraviolet-shape}
Under the hypotheses of Theorem~\ref{fp:thm-two-point} or
Corollary~\ref{fp:cor-cofinal-reference}, there are
\(c,C>0\) such that, for all sufficiently small \(d\) and all
admitted cutoffs, states and pairs,
\begin{equation}\label{fp:eq-ultraviolet-shape}
 \frac{c\vartheta^2 M_fM_h}{[\log(1/d)]^2}
 \le S^T_{2;a,\rho}(f,h)
 \le\frac{C}{[\log(1/d)]^2}
          \min\{M_fM_h,d^{-8}m_fm_h\}.
\end{equation}
For each fixed separated pair, any limit of the connected
correlations along admitted cutoffs obeys the same bounds.
If such limits exist for every pair in the Euclidean class and
\(\vartheta\le\vartheta_*\), the limit satisfies the ultraviolet bound on
\(\mathfrak T_{\mathbb R^4}(d;\vartheta_*)\) with
\(u(d)=C[\log(1/d)]^{-2}\). For the mass-one dilates constructed
below and \(\vartheta\le\vartheta_*\), the connected correlations,
and any corresponding fixed-pair limits, are bounded above and
below by positive constant multiples of
\(d^{-8}[\log(1/d)]^{-2}\).
On a family where the theorem gives
uniform vanishing of the relative error as \(d\downarrow0\),
the corresponding limiting correlations have leading coefficient
\(C_G/c_*^2\) multiplying \(\mathcal K(f,h)[\log(1/d)]^{-2}\).
\end{corollary}

\begin{proof}
At the witness, \(\zeta_K(d)\to1\) and
\(g_K(d)^{-2}/\log(1/d)\to c_*>0\) uniformly in the admitted
cutoff data. Shrink the distance threshold until the source
factor and the inverse-coupling ratio have positive lower
and finite upper bounds. The relative error bound
\(|r|\le1/2\) puts \(1+r\) in \([1/2,3/2]\).
Use the two-sided integral estimate
\eqref{fp:eq-shape-kernel-bounds} in the comparison formula of
Theorem~\ref{fp:thm-two-point}. This proves
\eqref{fp:eq-ultraviolet-shape}. Its constants are uniform
along the admitted cutoff sequence, so the inequalities
hold in any indicated limit. The inclusion at
\(\vartheta\le\vartheta_*\) follows from \eqref{fp:eq-shape}.

For an explicit pair at scale one, take
\(f(x)=(12/\pi^2)(1-|x|^2)_+^2\) and \(h(x)=f(x-4e_1)\).
The radial integral of \((1-|x|^2)_+^2\) is \(\pi^2/12\),
and its largest radial derivative has magnitude \(8/(3\sqrt3)\).
The support gap is two. These facts give
\((f,h)\in\mathfrak T_{\mathbb R^4}(1;\vartheta_*)\).
For a fixed \(x_0\), use
\(f_d(x)=d^{-4}f((x-x_0)/d)\) and
\(h_d(x)=d^{-4}h((x-x_0)/d)\).
Both masses stay one, and the amplitude product is a fixed
positive constant times \(d^{-8}\). Dilation and positive
amplitude scaling preserve the shape conditions at scale \(d\).
For these tests \eqref{fp:eq-ultraviolet-shape} gives the
claimed positive lower and upper bounds.

For the last assertion, divide
the comparison by \(C_G\mathcal K(f,h)/[c_*^2\log^2(1/d)]\).
The resulting three factors tend uniformly to one in the
order stated, and every fixed-pair cutoff limit inherits
the same uniform small-distance error.
\end{proof}

\subsection{Explicit tests at the fixed shape parameter}
The following check gives the explicit tests of the continuum application.
Let \(A_0\) be one on \([9/4,11/4]\), zero outside \([17/8,23/8]\),
and use the ramp \(3u^2-2u^3\), \(0\le u\le1\), on each transition
interval, with \(u\) increasing towards the plateau. Then
\(\int A_0=5/8\) and \(\|A_0'\|_\infty\le12\), by the integral
and maximum derivative of the ramp. For \(I=[-1/4,1/4]\) put
\(B_0=12\mathbf1_I^{*4}\). Direct integration of the triangular
convolution \(\mathbf1_I^{*2}\) gives
\[
 B_0(s)=
 \begin{cases}
  1-6s^2+6|s|^3,& |s|\le1/2,\\
  2(1-|s|)^3,&1/2\le|s|\le1,\\
  0,&|s|\ge1.
 \end{cases}
\]
It is \(C^2\), \(\|B_0\|_\infty=1\), and
\(\|B_0'\|_\infty=2\), with the maximum derivative attained at
\(|s|=1/3\). The Fourier transform is
\(\widehat B_0(q)=\tfrac34(\sin(q/4)/(q/4))^4\), with value
\(3/4\) at zero, by taking the fourth power of the transform of
\(\mathbf1_I\).

For transverse coordinates \(z\in\mathbb R^2\), take
\(q(s)=B_0(4s/3)\), \(D(s)=\sum_{j\in\mathbb Z}q(s-j)\) and
\(\chi=q/D\). The nearest integer yields \(D\ge2/27\), and at
most two summands are nonzero at each point. Since \(q\le D\),
the quotient rule gives \(\|\chi'\|_\infty\le108\).
The integer translates of \(\psi(z)=\chi(z_1)\chi(z_2)\) sum
to one by periodicity of \(D\). Let
\(\widetilde\psi(z)=\chi_0(z_1)\chi_0(z_2)\), where \(\chi_0=1\)
on \([-1,1]\), has support \([-3/2,3/2]\) and uses the same
ramp. The support of \(\psi\) lies in \([-3/4,3/4]^2\), so
\(0\le\psi\le\widetilde\psi\). The product rule gives
\(\|\nabla\psi\|_\infty\le216\) and
\(\|\nabla\widetilde\psi\|_\infty\le6\).
Integration gives \(\int\widetilde\psi=(5/2)^2\).

For a nonzero real \(\varphi\in C_c^1(\mathbb R^2)\),
\(\delta>0\), \(1\le\lambda\le2\) and \(i\in\mathbb Z^2\), set
\[
 A=\|\varphi\|_\infty+\delta\|\nabla\varphi\|_\infty,\qquad
 \psi_i(z)=\psi(z/\delta-i),\qquad
 \widetilde\psi_i(z)=\widetilde\psi(z/\delta-i).
\]
Put \(\varphi_i^-=2A\widetilde\psi_i\) and
\(\varphi_i^+=2A\widetilde\psi_i+\psi_i\varphi\).
The two functions
\[
 H_i^\pm(t,x,z)=A_0(t/\delta)B_0(4x/(\lambda\delta))
 \varphi_i^\pm(z)
\]
lie between \(A\) and \(3A\) times the product of the two
one-dimensional factors with \(\widetilde\psi_i\).
In particular their sup norms lie in \([A,3A]\).
The sum of \(H_i^+-H_i^-\) over the finitely many indices
with \(\psi_i\varphi\ne0\) equals
\(A_0(t/\delta)B_0(4x/(\lambda\delta))\varphi(z)\).
Their supports lie in balls of radius \(d=5\delta/2\), since
the sum of squared coordinate half-widths around
\((5\delta/2,0,\delta i)\) is at most \(5\delta^2\).
Each function and its time reflection have separation
between \(4\delta\) and \(6\delta\), by the time support
and the reflected displacement of a point on the plateau.
The mass is at least \(375A\delta^4/512\) by integrating the
product lower bound and using \(\lambda\ge1\).
The Lipschitz constant is bounded by
\(A[3(12+8)+2\cdot6+216+1]/\delta\).
Dividing the mass by \(3A(5\delta/2)^4\) gives \(1/160\).
The derivative bound is at most \(12000A/d\).
Thus \((H_i^\pm\circ\theta,H_i^\pm)\) belongs to
\(\mathfrak T_{\mathbb R^4}(5\delta/2;\vartheta_*)\),
where \(\theta(t,x,z)=(-t,x,z)\).


\begin{thebibliography}{Bal87}
\small
\bibitem{LP} \emph{Periodic Wilson states and reflection positivity}, companion manuscript, 2026.
\bibitem{MB} \emph{Linear lattice PDE with block constraints}, companion manuscript, 2026.
\bibitem{GI} \emph{Gaussian block integration and free action-density correlations}, companion manuscript, 2026.
\bibitem{RT} \emph{Coupling trajectories and partition stability in lattice Yang--Mills theory}, companion manuscript, 2026.
\bibitem{SI} \emph{Scalar sources and uniform insertions for lattice Yang--Mills fields}, companion manuscript, 2026.
\bibitem{Bal87} T.~Ba{\l}aban, \emph{Renormalization group approach to lattice gauge field theories. I. Generation of effective actions in a small field approximation and a coupling constant renormalization in four dimensions}, Comm. Math. Phys. \textbf{109} (1987), 249--301. \href{https://doi.org/10.1007/BF01215223}{doi:10.1007/BF01215223}.
\end{thebibliography}
\end{document}
